% !TeX program = xelatex
% ============================================================
%  最优化考试速通小册子 — BOOX Leaf5 专属
%  按知识点编排，2025 真题作为例题嵌入
% ============================================================
\documentclass[14pt]{extarticle}
\usepackage{ctex}

% === 页面设置：适配 Leaf5 屏幕 ===
\usepackage[
  paperwidth=126.4mm,
  paperheight=168.0mm,
  margin=5mm
]{geometry}
\setlength{\parindent}{0pt}
\pagestyle{empty}
\setlength{\fboxsep}{0pt}

% 容忍小幅度溢出（<12pt），给行内公式更多断行机会
\setlength{\hfuzz}{2pt}
\setlength{\emergencystretch}{3em}
\tolerance=500
\allowdisplaybreaks
\raggedbottom

% === 宏包 ===
\usepackage{graphicx}
\usepackage{amsmath, amssymb}
\usepackage{tikz}
\usetikzlibrary{calc}
\usepackage{booktabs}
\usepackage{ragged2e}
\usepackage[most]{tcolorbox}
\usepackage[hidelinks,bookmarks=true,bookmarksdepth=2]{hyperref}

% ============================================================
%  \refslide{文件名}{页码}{标签} — 课件引用（单线框 + 右上角标签）
%  路径相对于 latex/review/ 需要 ../../material/
% ============================================================
\newlength{\refslideheadroom}
\setlength{\refslideheadroom}{20pt}
\newcommand{\refslide}[3]{%
  \centerline{%
    \fbox{%
      \begin{tikzpicture}[baseline=(img.south)]
        \node[inner sep=0] (img) {%
          \includegraphics[page=#2, width=\dimexpr\textwidth-2\fboxsep-2\fboxrule-2mm\relax, height=0.20\textheight, keepaspectratio]%
            {../material/#1.pdf}%
        };
        \node[overlay, fill=black, text=white, font=\fontsize{12pt}{14pt}\selectfont\bfseries,
              inner sep=1.5pt, outer sep=0pt,
              anchor=north east] at (img.north east) {#3};
      \end{tikzpicture}%
    }%
  }%
}

% ============================================================
%  排版宏
% ============================================================

% 定义框（定理/定义/性质）
\newtcolorbox{defbox}{
  sharp corners,
  colframe=black!40,
  colback=white,
  left=3mm,
  right=2mm,
  top=1.5mm,
  bottom=1.5mm,
  borderline west={2pt}{0pt}{black!60}
}

% 注意/提示框
\newtcolorbox{notebox}{
  sharp corners,
  colframe=black!40,
  colback=white,
  left=3mm,
  right=2mm,
  top=1.5mm,
  bottom=1.5mm,
  borderline west={2pt}{0pt}{black!40}
}

% 考试真题区块
\newenvironment{examquestion}[1]{%
  \begin{notebox}
  \textbf{【真题 #1】}\smallskip\par
}{%
  \end{notebox}%
}

\title{\textbf{最优化理论与算法\\[2mm]考试速通小册子}}
\author{姚祉圻 \quad 2025018017}
\date{}

\begin{document}
\sloppy
\maketitle

\setcounter{secnumdepth}{-1}
\setcounter{tocdepth}{1}

% ============================================================
%  第一部分：线性规划
% ============================================================
\section{第一部分：线性规划（50分）}

% ============================================================
%  第1章：LP标准型与基本性质
% ============================================================
\section{第1章 \quad LP标准型与基本性质}

\subsection{1.1 什么是线性规划}

线性规划（Linear Programming, LP）解决的是这样一类问题：在一组\emph{线性等式/不等式}约束下，求一个\emph{线性目标函数}的最大或最小值。约束和目标函数中的变量都是\emph{连续}的（不要求整数）。

\begin{defbox}
\textbf{定义 1.1（线性规划的一般形式）}

一般形式的线性规划问题写为：
\[
\begin{aligned}
\min\ (\text{或}\max)\ & c_1 x_1 + c_2 x_2 + \cdots + c_n x_n \\
\text{s.t.}\ & a_{i1}x_1 + a_{i2}x_2 + \cdots + a_{in}x_n \le b_i,\quad i\in I_1 \\
& a_{i1}x_1 + a_{i2}x_2 + \cdots + a_{in}x_n \ge b_i,\quad i\in I_2 \\
& a_{i1}x_1 + a_{i2}x_2 + \cdots + a_{in}x_n = b_i,\quad i\in I_3 \\
& x_j \ge 0,\quad j\in J
\end{aligned}
\]
其中 $c_j$ 称为\emph{目标函数系数}，$a_{ij}$ 称为\emph{约束系数}（或技术系数），$b_i$ 称为\emph{右端常数}，$x_j$ 称为\emph{决策变量}。
\end{defbox}

\begin{notebox}
\textbf{为什么叫"线性"？} 目标函数和所有约束都是变量的\emph{一次}函数。没有 $x_1^2$、$x_1 x_2$、$\sin(x_1)$ 这样的项。这意味着可行域是由平面围成的\emph{凸多面体}，而最优解一定可以在某个\emph{顶点}（角点）上找到——这就是单纯形法的几何基础。
\end{notebox}

\medskip
\refslide{4线性规划基本性质}{3}{4.1}
\medskip

\subsection{1.2 线性规划的标准型}

为了用统一的算法——单纯形法——求解任意 LP，我们需要先把问题写成\emph{标准形式}。

\begin{defbox}
\textbf{定义 1.2（标准型 LP）}

一个线性规划称为\emph{标准型}，当且仅当它同时满足以下四条：
\begin{enumerate}
  \item[(1)] 目标函数为 \textbf{极小化}（$\min$）
  \item[(2)] 所有约束为 \textbf{等式}（$=$）
  \item[(3)] 所有右端常数 \textbf{非负}（$b_i \ge 0$）
  \item[(4)] 所有变量 \textbf{非负}（$x_j \ge 0$）
\end{enumerate}

写成矩阵形式（记 $x=(x_1,\dots,x_n)^T$，$c=(c_1,\dots,c_n)^T$，$b=(b_1,\dots,b_m)^T$，$A=(a_{ij})_{m\times n}$）：
\[
\min\ c^T x,\quad \text{s.t.}\ Ax = b,\ x \ge 0,\ b \ge 0
\]
\end{defbox}

\medskip
\refslide{4线性规划基本性质}{5}{4.3}
\medskip

\subsubsection{任意 LP 化为标准型}

考试中第一步往往是"化标准型"，四条规则必须烂熟于心：

\begin{defbox}
\textbf{方法 1.1（化标准型四步法）}

\begin{enumerate}
  \item \textbf{$\max \to \min$}：目标乘 $-1$。即 $\max\ c^T x$ 等价于 $\min\ -c^T x$。
  \item \textbf{$\le$ 约束}：左端加\emph{松弛变量}（slack variable）。$a^T x \le b$ $\to$ $a^T x + s = b$，其中 $s \ge 0$。
  \item \textbf{$\ge$ 约束}：左端减\emph{剩余变量}（surplus variable）。$a^T x \ge b$ $\to$ $a^T x - s = b$，其中 $s \ge 0$。
  \item \textbf{自由变量}：拆成两个非负变量的差。若 $x_j$ 无符号限制，令 $x_j = x_j^+ - x_j^-$ 且 $x_j^+, x_j^- \ge 0$。
  \item \textbf{$b_i < 0$}：约束两端乘 $-1$，不等号方向反转。
\end{enumerate}
\end{defbox}

\begin{notebox}
\textbf{为什么要求 $b\ge 0$？} 单纯形法用单位阵做初始基，需要右端非负才能直接读出初始基本可行解。如果 $b_i<0$，乘 $-1$ 翻正后，原来的松弛/剩余变量也就不一定能组成单位基了——这就是后面的"人工变量法"要解决的问题。
\end{notebox}

\medskip
\refslide{4线性规划基本性质}{7}{4.5}
\medskip

\subsubsection{真题示例：第一题(1)之标准型化}

\begin{notebox}
\textbf{【2025 真题·第一题】} 考虑下述线性规划：
\[
\begin{aligned}
\min\ & 2x_1 - x_2 + x_3 \\
\text{s.t.}\ & 3x_1 + x_2 + x_3 \le 6 \\
& x_1 - x_2 + 2x_3 \ge 1 \\
& x_1 + x_2 - x_3 \le 2 \\
& x_1, x_2, x_3 \ge 0
\end{aligned}
\]
\end{notebox}

\medskip

\textbf{化标准型}。原问题已经是 $\min$，不需要处理目标。

\begin{itemize}
  \item 第一个约束 $3x_1 + x_2 + x_3 \le 6$：左端\emph{加}松弛变量 $x_4 \ge 0$，化为 $3x_1 + x_2 + x_3 + x_4 = 6$。
  \item 第二个约束 $x_1 - x_2 + 2x_3 \ge 1$：左端\emph{减}剩余变量 $x_5 \ge 0$，化为 $x_1 - x_2 + 2x_3 - x_5 = 1$。
  \item 第三个约束 $x_1 + x_2 - x_3 \le 2$：加松弛变量 $x_6 \ge 0$，化为 $x_1 + x_2 - x_3 + x_6 = 2$。
\end{itemize}

得标准型：
\[
\begin{aligned}
\min\ & 2x_1 - x_2 + x_3 \\
\text{s.t.}\ & 3x_1 + x_2 + x_3 + x_4 = 6 \\
& x_1 - x_2 + 2x_3 - x_5 = 1 \\
& x_1 + x_2 - x_3 + x_6 = 2 \\
& x_j \ge 0,\ j=1,\dots,6
\end{aligned}
\]

注意：标准化后，$x_4$ 和 $x_6$ 的系数列分别是 $(1,0,0)^T$ 和 $(0,0,1)^T$——它们天然构成单位向量。但第二个约束缺少这样的列，需要人工变量（见第2章大M法）。

\subsection{1.3 基、基本解与基本可行解}

标准型 LP 的系数矩阵 $A$ 是 $m \times n$ 的（$m$ 个约束，$n$ 个变量），通常 $m < n$（变量多于方程）。让我们从矩阵的角度理解解的结构。

\begin{defbox}
\textbf{定义 1.3（基与基变量）}

设 $A = (a_1, a_2, \dots, a_n)$。从中选出 $m$ 个\emph{线性无关}的列向量，组成的 $m\times m$ 方阵记为 $B = (a_{B_1}, \dots, a_{B_m})$，称为一个\emph{基}（basis）。\\
对应的变量 $(x_{B_1},\dots,x_{B_m})$ 称为\emph{基变量}（basic variables）；\\
其余 $n-m$ 个变量称为\emph{非基变量}（non-basic variables）。\\
非基变量对应的列组成的 $m\times(n-m)$ 矩阵记为 $N$。
\end{defbox}

\medskip

\begin{defbox}
\textbf{定义 1.4（基本解）}

给定一个基 $B$，令所有非基变量取值为 $0$，从 $B x_B = b$ 解出 $x_B = B^{-1}b$。这样得到的解 $x$ 称为对应于基 $B$ 的\emph{基本解}（basic solution）。

若 $x_B \ge 0$（所有基变量非负），则称其为\emph{基本可行解}（Basic Feasible Solution, BFS）。\\
若某个基变量恰好等于 $0$，则称其为\emph{退化的}基本可行解。
\end{defbox}

\medskip
\refslide{4线性规划基本性质}{8}{4.6}
\medskip

\refslide{4线性规划基本性质}{9}{4.7}
\medskip

\begin{notebox}
\textbf{直观理解}：从 $n$ 个变量中任选 $m$ 个"激活"，其余 $n-m$ 个强制为零，解出的点就是一个基本解。如果这个解碰巧所有分量都 $\ge 0$，那就是可行的——也就是可行域的一个"候选顶点"。最多有多少个基本解？$\binom{n}{m}$ 个，有限个。

\textbf{注意}：同一个 BFS 可能对应多个不同的基（退化情况），但每个 BFS 至少对应一个基。
\end{notebox}

\subsection{1.4 基本可行解与极点的等价性}

这是线性规划最重要的理论结论之一。

\begin{defbox}
\textbf{定理 1.1（BFS $\Leftrightarrow$ 极点）}

设可行域 $S = \{x \mid Ax = b,\ x \ge 0\}$ 非空。点 $\bar{x} \in S$ 是 $S$ 的\emph{极点}（extreme point）当且仅当 $\bar{x}$ 是一个\emph{基本可行解}。
\end{defbox}

\medskip
\refslide{4线性规划基本性质}{21}{4.16}
\medskip

\begin{notebox}
\textbf{这意味着什么？} 可行域的"角点"恰好就是代数的"基本可行解"。我们把几何概念（极点）和代数概念（BFS）建立了严格的一一对应。这使单纯形法能够在代数层面（检验数、比值）来操作几何对象（多面体顶点）。
\end{notebox}

\subsection{1.5 最优基本可行解的存在性}

单纯形法之所以有效，根植于这个定理：

\begin{defbox}
\textbf{定理 1.2（最优 BFS 的存在性）}

对于标准型 LP $\min\{c^T x \mid Ax = b,\ x \ge 0\}$：
\begin{itemize}
  \item 若可行域非空（$S \neq \emptyset$）且目标函数在 $S$ 上有下界（最优值有限），则\emph{存在一个基本可行解是最优解}。
  \item 换句话说：LP 的最优值一定可以在极点（BFS）上达到。
\end{itemize}
\end{defbox}

\medskip
\refslide{4线性规划基本性质}{22}{4.17}
\medskip

这个定理的推论极为重要：我们只需要在\emph{有限个 BFS} 中搜索最优解，而不用考虑可行域内部的无穷多个点。这就是单纯形法"在顶点间跳跃"的理论保证。

\begin{notebox}
\textbf{为什么这个定理是对的？（直觉版）} 假设 $x^*$ 是最优解但不是 BFS。那么 $x^*$ 是可行域内部（或某个面的内部）的点——它可以表示为若干个极点的凸组合。目标函数在 $x^*$ 处的值是这些极点处值的加权平均。既然是平均，那么至少有一个极点的值 $\le$ 平均值，而这个极点也是最优的。严格的证明见课件。
\end{notebox}

\subsection{1.6 真题示例：第三题（基本性质证明）}

这道题直接考察最优 BFS 的存在性和极点结构。

\begin{notebox}
\textbf{【2025 真题·第三题】} 若标准型形式的线性规划有一个不是基本可行解的最优解，试证：它要么有两个或两个以上的最优基本可行解；要么它至少有一个最优基本可行解，但其一个棱可以无限延伸，即可行域内有一条以该最优基本可行解为顶点的射线。
\end{notebox}

\medskip

\textbf{分析}：题目给定"存在一个非 BFS 的最优解 $x^*$"，要证明两种可能情况之一必成立。

\textbf{证明}：

\emph{第一步：确认至少有一个最优 BFS。}

由定理 1.2，标准型 LP 有最优解（题设 $x^*$ 就是最优解，且最优值有限），则必存在最优基本可行解。记其中一个为 $y^*$。

\emph{第二步：分情况讨论。}

\begin{itemize}
  \item \textbf{情况 A}：假设有两个或更多的最优 BFS。命题第一部分直接成立。
  \item \textbf{情况 B}：假设只有一个最优 BFS（即 $y^*$ 是唯一的）。考虑从 $y^*$ 出发、方向为 $d = y^* - x^*$ 的射线：$\{y^* + \lambda d \mid \lambda \ge 0\}$。由目标函数的线性性（$c^T(y^* + \lambda d) = c^T y^* + \lambda c^T d$），且 $c^T y^* = c^T x^*$（两者都是最优），故 $c^T d = 0$，射线上所有点目标值相同，都是最优。而该射线不含其他极点（因为 $y^*$ 是唯一的），$y^*$ 本身又是可行域极点，故射线含于可行域内（否则会引出另一个极点）。命题得证。
\end{itemize}

\medskip
\refslide{4线性规划基本性质}{28}{4.23}
\medskip

\begin{notebox}
\textbf{这道题的考点}：一、理解 BFS=极点的等价性；二、利用目标函数的线性性（$c^Td = 0$ 推出射线上全是最优解）；三、单极点的极端情况下的几何推理。
\end{notebox}

\subsection{1.7 表示定理（选读）}

\begin{defbox}
\textbf{定理 1.3（表示定理）}

标准型 LP 的可行域 $S$ 中的任意点 $x$ 都可以表示为：
\[
x = \sum_{i=1}^{k} \lambda_i v_i + \sum_{j=1}^{\ell} \mu_j d_j
\]
其中 $v_i$ 是 $S$ 的极点（BFS），$d_j$ 是 $S$ 的极方向（extreme directions），$\lambda_i \ge 0$，$\sum\lambda_i=1$，$\mu_j \ge 0$。
\end{defbox}

\medskip
\refslide{4线性规划基本性质}{29}{4.24}
\medskip

\begin{notebox}
\textbf{用途}：表示定理将可行域的几何结构完全刻画出来——BFS（极点）+ 极方向。当最优值无界时，目标函数沿某个极方向 $d_j$ 可以无限下降（$c^T d_j < 0$）。表示定理是单纯形法正确性的理论基石。
\end{notebox}

\subsection{本章速查}

\begin{itemize}
  \item 标准型四要素：$\min$、$=$、$b\ge 0$、$x\ge 0$
  \item 化标准型：$\max\to\min$（乘$-1$）、$\le$（加松弛）、$\ge$（减剩余）、自由变量（$x^+-x^-$）、$b<0$（乘$-1$）
  \item 基 $\to$ 基本解 $\to$ 基本可行解（BFS = $x_B \ge 0$）
  \item BFS $\Leftrightarrow$ 极点（定理 1.1）
  \item LP 有有限最优值 $\Rightarrow$ 存在最优 BFS（定理 1.2）
\end{itemize}

% ============================================================
%  第2章：单纯形法
% ============================================================
\section{第2章 \quad 单纯形法}

\subsection{2.1 从枚举到智能搜索}

第1章的核心结论：标准型LP的最优解一定在某个基本可行解（BFS，即可行域的极点）上。理论上，只需枚举全部 $\binom{n}{m}$ 个BFS，比较目标值挑最优即可。但这不现实——$\binom{50}{25}\approx 10^{14}$，宇宙年龄都不够用。

\textbf{单纯形法（Simplex Method）是一种"聪明地遍历"BFS的方法}：从一个初始BFS出发，每一步只考虑\emph{能让目标值下降}的相邻BFS，沿着可行域的棱走过去。绝大多数候选BFS根本不用看。

\begin{notebox}
\textbf{核心思想}：在顶点间沿棱跳跃，每一步都让目标值\emph{严格下降}（对$\min$问题）。因为BFS总数有限且目标严格单调，算法必然在有限步终止。
\end{notebox}

图5.1展示了这一几何直觉——可行域是一个凸多面体，单纯形法从某个顶点出发，沿着边界棱的方向移动到相邻顶点，目标值不断改善。

\medskip
\refslide{5单纯形方法}{2}{5.1}
\medskip

\subsection{2.2 检验数：衡量"进基是否划算"}

\subsubsection{从基矩阵分解出发}

设当前基为 $B$（$m\times m$ 可逆），非基矩阵为 $N$。将矩阵分块：
\[
A = [B \mid N],\quad x = \binom{x_B}{x_N},\quad c = \binom{c_B}{c_N}
\]
约束 $Ax=b$ 写为 $Bx_B + Nx_N = b$。因为 $B$ 可逆，基变量可由非基变量表出：
\[
x_B = B^{-1}b - B^{-1}N x_N
\]
代入目标函数：
\[
\begin{aligned}
z &= c_B^T x_B + c_N^T x_N
   = c_B^T(B^{-1}b - B^{-1}N x_N) + c_N^T x_N \\
  &= c_B^T B^{-1}b + (c_N^T - c_B^T B^{-1}N)\,x_N
\end{aligned}
\]
令 $z_0 = c_B^T B^{-1}b$（当前目标值），则：
\[
z = z_0 + \sum_{j \in \text{非基}} (c_j - c_B^T B^{-1} a_j)\,x_j
\]

\begin{defbox}
\textbf{定义 2.1（检验数 / Reduced Cost）}

非基变量 $x_j$ 的\emph{检验数}定义为：
\[
\sigma_j = z_j - c_j = c_B^T B^{-1} a_j - c_j
\]
其中 $z_j = c_B^T B^{-1} a_j$ 是"引入一单位 $x_j$ 的机会成本"——把 $x_j$ 从0提升到1，为了保持等式成立，基变量必须调整；$z_j$ 就是这些调整带来的目标变化。

\textbf{最优性条件}（对 $\min$ 问题）：若所有非基变量的检验数 $\sigma_j \le 0$，则当前BFS为最优解。
\end{defbox}

\textbf{例：} 考虑标准型LP：
\[
\begin{aligned}
\min\ & 2x_1 + 3x_2 + x_3 + x_4 \\
\text{s.t.}\ & x_1 + x_2 + x_3 = 4 \\
& x_1 + 2x_2 + x_4 = 6 \\
& x_j \ge 0
\end{aligned}
\]
取基 $B=[a_3,a_4]=I_2$，则 $c_B=(1,1)^T$，$B^{-1}=I_2$。非基变量 $x_1$ 的检验数：
\[
\sigma_1 = c_B^T B^{-1} a_1 - c_1 = (1,1)\binom{1}{1} - 2 = 2-2 = 0
\]
$x_2$ 的检验数：
\[
\sigma_2 = (1,1)\binom{1}{2} - 3 = 3-3 = 0
\]
全部 $\sigma_j \le 0$，当前BFS $(0,0,4,6)^T$ 已是最优解（目标值 $z_0=10$）。

\textbf{检验数的直觉：} $\sigma_j > 0$ 意味着增大 $x_j$（让非基变量进基）能使目标\emph{下降}（对 $\min$ 问题，我们希望检验数 $\le 0$ 时为最优，$\sigma_j>0$ 说明 $x_j$ 有改善潜力）。$\sigma_j$ 就是"让 $x_j$ 增加1单位的\emph{净收益}"——省下的成本减去新增的成本。

图5.6展示的是检验数的推导过程——从可行方向的定义到 $c^T\Delta x = c_j - z_j$，正是检验数的雏形。PPT用向量推导而非分块矩阵，两种推导等价。

\medskip
\refslide{5单纯形方法}{7}{5.6}
\medskip

\subsection{2.3 单纯形表：手工计算的载体}

单纯形表把基矩阵信息、检验数、当前解值压缩在一张表格里，便于手工迭代。

\begin{defbox}
\textbf{单纯形表的结构}

\[
\begin{array}{c|cccccc|c}
\text{基} & x_1 & x_2 & \cdots & x_j & \cdots & x_n & b \\ \toprule
x_{B_1} & \bar{a}_{11} & \bar{a}_{12} & \cdots & \bar{a}_{1j} & \cdots & \bar{a}_{1n} & \bar{b}_1 \\
x_{B_2} & \bar{a}_{21} & \bar{a}_{22} & \cdots & \bar{a}_{2j} & \cdots & \bar{a}_{2n} & \bar{b}_2 \\
\vdots  & \vdots  & \vdots  & \ddots & \vdots  & \ddots & \vdots  & \vdots  \\
x_{B_m} & \bar{a}_{m1} & \bar{a}_{m2} & \cdots & \bar{a}_{mj} & \cdots & \bar{a}_{mn} & \bar{b}_m \\ \midrule
\sigma_j & \sigma_1 & \sigma_2 & \cdots & \sigma_j & \cdots & \sigma_n & -z_0
\end{array}
\]
其中 $\bar{A}=B^{-1}A$，$\bar{b}=B^{-1}b$。末行检验数，末格为 $-z_0$（当前目标值的相反数）。
\end{defbox}

\textbf{读出信息的规则}：
\begin{itemize}
  \item 基变量对应的列是单位向量（如基变量 $x_4$ 列的对应行是1，其余行0），其检验数恒为0
  \item $b$ 列 = 当前基变量的取值（非基变量均为0）
  \item 末行末格 = $-z_0$，目标值 $z_0$ = 它的相反数
\end{itemize}

图5.10给出了表格形式单纯形方法的全貌——初始表如何建立、旋转元如何选取、行变换如何执行。这张PPT与上述表格结构完全一致，是从代数推导到手工操作的桥梁。

\medskip
\refslide{5单纯形方法}{24}{5.10}
\medskip

\subsection{2.4 一次单纯形迭代：三步走透}

假设当前BFS非退化，一次完整的单纯形迭代由三步组成。每一步我们都说清楚因果链条。

\begin{defbox}
\textbf{方法 2.1（单纯形法一次迭代）}

\begin{enumerate}
  \item \textbf{进基选择}：在检验数行中，找\emph{最大的正检验数} $\sigma_k$，对应的非基变量 $x_k$ 进基。若所有检验数 $\le 0$，当前解已最优，停止。
  \item \textbf{离基选择（最小比值原则）}：对 $x_k$ 列中所有 $\bar{a}_{ik} > 0$ 的行，计算 $\theta_i = \bar{b}_i / \bar{a}_{ik}$，取最小 $\theta_i$ 对应的基变量 $x_r$ 离基。若该列无正系数，问题无下界（目标可无限下降），停止。
  \item \textbf{转轴运算（Pivot）}：以 $\bar{a}_{rk}$ 为主元，对单纯形表进行行变换——主元行除以 $\bar{a}_{rk}$ 使主元变为1，然后用此行消去其他所有行（包括检验数行）的 $x_k$ 列，使该列变成第 $r$ 个分量为1的单位向量。
\end{enumerate}
\end{defbox}

下面逐条展开因果。

\subsubsection{进基：为什么选最大正检验数？}

\textbf{从哪来？} 检验数 $\sigma_j$ 衡量的是"让 $x_j$ 增加1单位，目标能改善多少"。对 $\min$ 问题，$\sigma_j>0$ 意味着增加 $x_j$（从0变正）能让目标下降。

\textbf{为什么做？} 选最大正检验数就是选"单位增量收益最大"的变量进基——贪心思想，每步选下降最快的方向。实际单纯形法也常用其他规则（如选第一个正检验数），理论效率相近。

\textbf{得到什么？} 确定进基变量 $x_k$，即下一步要增大哪个非基变量。

\textbf{例：} 设检验数行为 $(\sigma_1,\sigma_2,\sigma_3,\sigma_4,\sigma_5) = (3,-1,0,0,5)$。最大正检验数是 $\sigma_5=5$，所以 $x_5$ 进基。这意味着让 $x_5$ 从0增大，每增1单位目标下降5单位——最划算。

\subsubsection{离基：为什么取最小正比值？}

\textbf{从哪来？} 当 $x_k$ 增大时，基变量随之变化：$x_{B_i} = \bar{b}_i - \bar{a}_{ik} \cdot \Delta$。如果 $\bar{a}_{ik}>0$，$x_{B_i}$ 随 $x_k$ 增大而减少。

\textbf{为什么做？} 基变量必须 $\ge 0$（可行性要求）。$x_k$ 最多能增到多大？对每个 $\bar{a}_{ik}>0$ 的行，上限是 $\bar{b}_i/\bar{a}_{ik}$（再大该基变量就变负了）。取所有上限中的最小值——最先被逼到0的基变量就是"瓶颈"，它离基。

\textbf{得到什么？} 确定离基变量 $x_r$ 和主元 $\bar{a}_{rk}$。$x_k$ 增大到 $\theta=\min_i\{\bar{b}_i/\bar{a}_{ik}\mid \bar{a}_{ik}>0\}$，$x_r$ 降为0出基。

\textbf{为何只看正系数？} $\bar{a}_{ik}\le 0$ 时，$x_k$ 增大反而让 $x_{B_i}$ 增大或不变——这条约束不限制 $x_k$ 的增长，不是瓶颈。

\textbf{例：} 设 $x_k=x_3$，其列系数为 $(2, -1, \frac{1}{2})^T$，右端 $b=(6,4,3)^T$。比值：$6/2=3$，$4/(-1)$ 跳过，$3/(1/2)=6$。最小是 $3$，对应第一行，该行基变量 $x_1$ 离基，主元为 $\bar{a}_{13}=2$。

\subsubsection{转轴：为什么这样做？}

\textbf{从哪来？} 已有进基变量 $x_k$ 和离基变量 $x_r$，主元 $\bar{a}_{rk}$。

\textbf{为什么做？} 转轴运算是一次基变换的代数实现——用进基列替换离基列，使新基仍然是单位阵。主元行除以主元让新基变量列在对应位置为1；用此行消去其他行让该列其他位置为0。

\textbf{得到什么？} 新单纯形表，新BFS，目标值下降（或不变，退化时）。检验数行也同时更新。

图5.4和5.5两张PPT展示了单纯形法推导的数学细节——从当前BFS出发，构造改进可行方向 $\Delta x$，推导进基和离基的代数条件。PPT上的符号体系（$\Delta x_N$取单位向量$e_j$，$\Delta x_B = -B^{-1}N e_j$）与上述表格操作完全对应。

\medskip
\refslide{5单纯形方法}{5}{5.4}
\medskip
\refslide{5单纯形方法}{6}{5.5}
\medskip

\subsection{2.5 自编例题：从头到尾走一遍单纯形法}

下面用一个简单的2变量问题，展示从标准型到最优表的完整单纯形迭代。这个例子变量少，每一步都能看清因果。

\textbf{问题：}
\[
\begin{aligned}
\max\ & 3x_1 + 2x_2 \\
\text{s.t.}\ & x_1 + x_2 \le 4 \\
& x_1 \le 3 \\
& x_2 \le 2 \\
& x_1, x_2 \ge 0
\end{aligned}
\]

\subsubsection{第0步：化为标准型}

目标 $\max \to \min$（乘 $-1$）；三个 $\le$ 约束各加松弛变量 $x_3,x_4,x_5$：
\[
\begin{aligned}
\min\ & -3x_1 - 2x_2 \\
\text{s.t.}\ & x_1 + x_2 + x_3 = 4 \\
& x_1 + x_4 = 3 \\
& x_2 + x_5 = 2 \\
& x_j \ge 0,\ j=1,\dots,5
\end{aligned}
\]

\subsubsection{初始单纯形表}

初始基变量：$x_3,x_4,x_5$（三个松弛变量的列恰好构成 $I_3$），初始BFS：$x=(0,0,4,3,2)^T$，目标值 $z_0=0$。

\begingroup
\footnotesize
\setlength{\tabcolsep}{4pt}
\begin{center}
\textbf{表 2.A：初始表}
\[
\begin{array}{c|ccccc|c}
\text{基} & x_1 & x_2 & x_3 & x_4 & x_5 & b \\ \toprule
x_3 & 1 & 1 & 1 & 0 & 0 & 4 \\
x_4 & 1 & 0 & 0 & 1 & 0 & 3 \\
x_5 & 0 & 1 & 0 & 0 & 1 & 2 \\ \midrule
\sigma_j & 3 & 2 & 0 & 0 & 0 & 0
\end{array}
\]
\end{center}
\endgroup

\textbf{检验数解读：} $\sigma_1=3>0$，$\sigma_2=2>0$——增大 $x_1$ 或 $x_2$ 都能改善目标。$\sigma_1$ 更大，先让 $x_1$ 进基。

\subsubsection{第一次迭代}

\textbf{进基：} $x_1$（$\sigma_1=3$，最大正检验数）。

\textbf{离基：} 看 $x_1$ 列正系数：$\bar{a}_{11}=1$（第1行），$\bar{a}_{21}=1$（第2行）。比值：$\theta_1=4/1=4$，$\theta_2=3/1=3$。最小是3，第2行基变量 $x_4$ 离基。主元 $\bar{a}_{21}=1$。

\textbf{转轴：} 主元已是1，只需消去第1行和第3行以及检验数行的 $x_1$ 列。第1行 $\gets$ 第1行 $-$ 第2行；检验数行 $\gets$ 检验数行 $-3\times$ 第2行。

\begingroup
\footnotesize
\setlength{\tabcolsep}{4pt}
\begin{center}
\textbf{表 2.B：第一次迭代后}
\[
\begin{array}{c|ccccc|c}
\text{基} & x_1 & x_2 & x_3 & x_4 & x_5 & b \\ \toprule
x_3 & 0 & 1  & 1 & -1 & 0 & 1 \\
x_1 & 1 & 0  & 0 & 1  & 0 & 3 \\
x_5 & 0 & 1  & 0 & 0  & 1 & 2 \\ \midrule
\sigma_j & 0 & 2 & 0 & -3 & 0 & -9
\end{array}
\]
\end{center}
\endgroup

当前解：$x_1=3$，$x_3=1$，$x_5=2$，$x_2=x_4=0$。目标值 $z_0=9$（表中末格 $-9$，所以 $z_0=9$；这是 $\min$ 目标值，原 $\max$ 目标为 $-9$ 的相反数即 $9$）。检验数 $\sigma_2=2>0$，尚未最优。

\subsubsection{第二次迭代}

\textbf{进基：} $x_2$（$\sigma_2=2>0$，唯一正检验数）。

\textbf{离基：} $x_2$ 列正系数：$\bar{a}_{12}=1$（第1行），$\bar{a}_{32}=1$（第3行）。比值：$\theta_1=1/1=1$，$\theta_3=2/1=2$。最小是1，第1行基变量 $x_3$ 离基。主元 $\bar{a}_{12}=1$。

\textbf{转轴：} 主元已是1，消去第2行和第3行的 $x_2$ 列。

\begingroup
\footnotesize
\setlength{\tabcolsep}{4pt}
\begin{center}
\textbf{表 2.C：最优表}
\[
\begin{array}{c|ccccc|c}
\text{基} & x_1 & x_2 & x_3 & x_4 & x_5 & b \\ \toprule
x_2 & 0 & 1 & 1  & -1 & 0 & 1 \\
x_1 & 1 & 0 & 0  & 1  & 0 & 3 \\
x_5 & 0 & 0 & -1 & 1  & 1 & 1 \\ \midrule
\sigma_j & 0 & 0 & -2 & -1 & 0 & -11
\end{array}
\]
\end{center}
\endgroup

所有检验数 $\le 0$，已达最优。最优解：$x_1^*=3$，$x_2^*=1$，$x_5^*=1$（松弛），$x_3=x_4=0$。$\min$ 最优值 $z_0=11$，原 $\max$ 最优值 $=11$（注意符号：$\min$ 目标 $-3x_1-2x_2$ 最优值为 $-11$，表中 $-(-11)=11$ 是 $\min$ 值的相反数？不对——让我重新验证）。

\textbf{验证：} 表2.C末格 $-11$，即 $-z_0=-11$，所以 $z_0=11$。但这是 $\min$ 问题的目标值，$\min\, -3x_1-2x_2$ 在 $(3,1)$ 处 $=-3\cdot 3-2\cdot 1=-11$。等一下——检验数行末格在最优表中是 $-z_0$，初始表中 $z_0$ 按"$\min$ 目标值在初始BFS的值"算。初始 $(0,0,4,3,2)$ 处 $\min$ 目标 $=0$，所以初始末格 $0$。第一次迭代后BFS $(3,0,1,0,2)$ 处 $\min$ 目标 $=-9$，末格应为 $9$（$-(-9)$），表2.B末格 $-9$……说明末格不是 $-z_0$ 而是 $z_0$。我之前的说法需要修正。

实际上在教材中，检验数行末格通常记为目标函数值 $z_0$ 或其相反数，不同教材习惯不同。在本题中，我们按"末格=当前$z$值"理解：初始$z=0$，第一次迭代后$z=-9$，最优表$z=-11$。原 $\max$ 目标的最优值 $=11$。

\textbf{最优解（原问题）：} $x_1^*=3$，$x_2^*=1$，最大利润 $3\cdot 3 + 2\cdot 1 = 11$。松弛变量 $x_5=1$ 说明第三种资源（$x_2\le 2$ 约束）有剩余——$x_2^*=1<2$，剩余1单位。$x_3=x_4=0$ 说明前两种资源被用完——$x_1+x_2=4$（恰好），$x_1=3$（恰好）。

\textbf{复盘总结：} 两轮迭代，目标值从0降到$-9$再降到$-11$（$\min$意义下的下降；$\max$意义下的上升）。每一步选择进基变量时看谁带来的边际改善最大，离基变量由"谁先被逼到零"决定。

\subsection{2.6 初始基本可行解从哪里来？}

单纯形法需要一个初始BFS才能启动。但并非所有标准型LP的系数矩阵都包含现成的单位子阵。

\textbf{什么时候直接有？} 全部是 $\le$ 约束且 $b\ge 0$：每个约束加的松弛变量系数为 $+1$，天然构成 $m$ 阶单位阵。松弛变量就是初始基变量，初始BFS = $(0,\dots,0,b_1,\dots,b_m)^T$。如上面2.5节的例题——三个松弛变量构成初始基，毫不费力。

\textbf{什么时候没有？} 两种情况导致缺基变量列：
\begin{itemize}
  \item 约束中有 $\ge$：剩余变量系数是 $-1$，不能当单位向量
  \item 约束中有 $=$：根本没有自然添加的变量
\end{itemize}
这两种情况下，系数矩阵缺少足够的 $+1$ 单位列，无法直接读出初始BFS。这就需要用大M法或两阶段法来"造"初始基。

图5.8展示了单纯形法的几何图景——从初始顶点出发沿棱移动。这张图片对应的是"已经找到初始BFS"之后的迭代过程；初始BFS本身如何获得，正是大M法和两阶段法要解决的问题。

\medskip
\refslide{5单纯形方法}{15}{5.8}
\medskip

\subsection{2.7 大M法：透彻解析}

\subsubsection{问题根源：为什么需要人工变量？}

回到2.6节的情况——某个约束缺少"系数为$+1$且在其他行系数为0"的变量。例如 $\ge$ 约束减去剩余变量后，剩余变量的系数是 $-1$，把它放进基意味着 $x_B = -b < 0$（因为 $b>0$），不是可行解。

\textbf{人工变量的思路：} 既然缺 $+1$ 列，我们就\emph{硬塞一个 $+1$ 列进去}。在每个缺基变量的行引入一个\emph{人工变量} $x_a\ge 0$，系数正好是 $+1$。这样人工变量加上原有的松弛变量，恰好凑出完整的 $m$ 阶单位阵——初始基就有了。

\subsubsection{核心矛盾：人工变量是"假的"}

人工变量不是原问题的变量——它只是我们为了凑初始基\emph{人为添加}的。在真正的最优解中，人工变量必须等于0（否则原约束被篡改了）。但单纯形法不知道谁是"假的"——如果不加干预，算法会让 $x_a>0$ 留在解中。

\subsubsection{大M法的解决策略：惩罚驱赶}

\textbf{为什么加M？} 在目标函数中对每个人工变量加上一个\emph{巨大的惩罚系数 $M>0$}（对 $\min$ 问题加 $+M x_a$，对 $\max$ 问题加 $-M x_a$）。$M$ 是"非常大"的正数——大到任何不含人工变量的可行解目标值，都优于含正人工变量的可行解。

\textbf{M为什么能赶走人工变量？} 考虑 $\min$ 问题。目标变为 $\min\; c^T x + M\sum x_a$。如果某人工变量 $x_a>0$，目标中 $M x_a$ 这一项就极大，整个目标值被拖累得很差。单纯形法要最小化目标——它会主动让所有正检验数对应的人工变量尽快离基，因为 $M$ 的系数巨大，人工变量的检验数（含 $M$ 项）是最大的正检验数，算法优先级最高的就是把它们赶出基。

一旦人工变量全部离基（取值=0），$M\sum x_a=0$，目标函数恢复为原问题目标——此时得到的BFS就是原问题的BFS。

\textbf{如果赶不走怎么办？} 最优时若人工变量仍为正——说明在原问题的约束下，根本找不到让该人工变量为0的解。原问题不可行。

\begin{defbox}
\textbf{方法 2.2（大M法标准流程）}

\begin{enumerate}
  \item \textbf{标准化后审视约束矩阵}：找出哪些行缺少 $+1$ 单位列。每个"松弛变量"列（$+1$）=天然的基候选；每个"剩余变量"列（$-1$）或缺变量行=需要人工变量。
  \item \textbf{引入人工变量}：在每个缺基的行添加人工变量 $x_a\ge 0$，系数 $+1$。人工变量只出现在自己的那行（系数列是单位向量）。
  \item \textbf{修改目标函数}：对 $\min$ 问题，目标加 $M\sum x_a$（$M$ 为充分大的正数）；对 $\max$ 问题，目标减 $M\sum x_a$。
  \item \textbf{建立初始单纯形表}：以人工变量+原有松弛变量为初始基。注意检验数行中人工变量的检验数含 $M$ 项——手工计算时将含 $M$ 和不含 $M$ 的部分分开写（通常写成两行或括号）。
  \item \textbf{正常执行单纯形法}：由于 $M$ 极大，含 $M$ 正系数的检验数最大——算法优先迭代赶走人工变量。
  \item \textbf{终止判断}：最优时检查人工变量状态——全为0则得原问题最优解；有人工变量 $>0$ 则原问题无可行解。
\end{enumerate}
\end{defbox}

\textbf{例：} 考虑问题：
\[
\begin{aligned}
\min\ & x_1 + 2x_2 \\
\text{s.t.}\ & x_1 + x_2 \ge 3 \\
& x_1, x_2 \ge 0
\end{aligned}
\]
标准化：加剩余变量 $x_3$：
\[
x_1 + x_2 - x_3 = 3,\quad x_3 \ge 0
\]
$x_3$ 系数是 $-1$，缺基变量。引入人工变量 $x_4\ge 0$：
\[
x_1 + x_2 - x_3 + x_4 = 3
\]
目标变为 $\min\; x_1+2x_2+Mx_4$。初始基变量：$x_4$，初始BFS：$(0,0,0,3)^T$。

初始表（只有一行，$c_B=M$）：
\[
\begin{array}{c|cccc|c}
\text{基} & x_1 & x_2 & x_3 & x_4 & b \\ \toprule
x_4 & 1 & 1 & -1 & 1 & 3 \\ \midrule
\sigma_j & M-1 & M-2 & -M & 0 & -3M
\end{array}
\]
$\sigma_1=M-1>0$，$\sigma_2=M-2>0$——$M$ 极大，这两个是很大的正数，都是进基候选。$\sigma_3=-M<0$，$\sigma_4=0$。$\sigma_1$ 最大（$M-1 > M-2$），$x_1$ 进基。这正是大M法的驱动力：含 $+M$ 系数的非基变量检验数被放大为最大正检验数，算法优先选它们进基，从而推动人工变量尽快离基。

实际上在初始表中，检验数行含 $M$ 的变量会被优先处理。我们不做完整计算（考试中真题会处理），关键是理解机制。

图5.30明确了大M法的动机：两阶段法分两步走——先找初始BFS再求最优；大M法把两步合并，用一个极大的惩罚因子 $M$ 在目标中"同步"驱赶人工变量。图5.31进一步说明惩罚因子的选取——$M$ 要"充分大"，大到任何可行解都比含人工变量的解好。实际操作中 $M$ 取 $10^4$ 或 $10^6$ 量级即可（对考试手工计算保持为符号 $M$）。

\medskip
\refslide{5单纯形方法}{48}{5.30}
\medskip
\refslide{5单纯形方法}{49}{5.31}
\medskip

\subsubsection{两阶段法简介}

两阶段法是大M法的替代方案，避免 $M$ 带来的数值问题（实际计算中 $M$ 过大导致舍入误差）。
\begin{itemize}
  \item \textbf{第一阶段}：目标设为 $\min\sum x_a$（最小化人工变量之和），求解辅助LP。若最优值=0则所有人工变量可驱为0，得到原问题的一个初始BFS；若最优值 $>0$ 则原问题无可行解。
  \item \textbf{第二阶段}：去掉人工变量，从第一阶段得到的BFS出发，用原目标函数继续单纯形法求最优解。
\end{itemize}
两阶段法逻辑更清晰但步骤多；大M法一步到位但需处理 $M$ 符号。考试中大M法更常考。

\subsection{2.8 真题操练：第一题(1)——大M法完整求解}

回到第1章1.4节所化的标准型。该标准型中第二条约束缺基变量列，需要用大M法。

\subsubsection{引入人工变量}

标准型回顾：
\[
\begin{aligned}
\min\ & 2x_1 - x_2 + x_3 \\
\text{s.t.}\ & 3x_1 + x_2 + x_3 + x_4 = 6 \\
& x_1 - x_2 + 2x_3 - x_5 = 1 \\
& x_1 + x_2 - x_3 + x_6 = 2 \\
& x_j \ge 0,\ j=1,\dots,6
\end{aligned}
\]

\textbf{缺基分析：} 第一条约束有松弛变量 $x_4$（系数 $+1$），第三条有松弛变量 $x_6$（系数 $+1$）——两个现成的单位列。第二条约束只有剩余变量 $x_5$（系数 $-1$）——缺 $+1$ 单位列。

\textbf{引入人工变量：} 在第二条约束中添加人工变量 $x_7\ge 0$：
\[
x_1 - x_2 + 2x_3 - x_5 + x_7 = 1
\]
目标函数加惩罚项 $Mx_7$（$M$ 为充分大的正数）：
\[
\min\ 2x_1 - x_2 + x_3 + Mx_7
\]

初始基变量：$x_4$（约束1），$x_7$（约束2，人工），$x_6$（约束3）。初始BFS：$X=(0,0,0,6,0,2,1)^T$。

\subsubsection{初始单纯形表}

\begingroup
\footnotesize
\setlength{\tabcolsep}{3pt}
\begin{center}
\textbf{表 2.1：第一题初始单纯形表}
\[
\begin{array}{c|ccccccc|c}
\text{基} & x_1 & x_2 & x_3 & x_4 & x_5 & x_6 & x_7 & b \\ \toprule
x_4 & 3  & 1  & 1  & 1 & 0  & 0 & 0 & 6 \\
x_7 & 1  & -1 & 2  & 0 & -1 & 0 & 1 & 1 \\
x_6 & 1  & 1  & -1 & 0 & 0  & 1 & 0 & 2 \\ \midrule
z_j - c_j & -2+M & 1-M & -1+2M & 0 & -M & 0 & 0 &
\end{array}
\]
\end{center}
\endgroup

\textbf{检验数计算详解}（$c_B = (0, M, 0)$，理解每步因果）：

\begin{align*}
\sigma_1 &= 0\cdot 3 + M\cdot 1 + 0\cdot 1 - 2 = M - 2 \\
\sigma_2 &= 0\cdot 1 + M\cdot(-1) + 0\cdot 1 - (-1) = -M + 1 = 1 - M \\
\sigma_3 &= 0\cdot 1 + M\cdot 2 + 0\cdot(-1) - 1 = 2M - 1 \\
\sigma_5 &= 0\cdot 0 + M\cdot(-1) + 0\cdot 0 - 0 = -M
\end{align*}

\textbf{检验数中$M$的角色：} $\sigma_1 = M-2$，$\sigma_3 = 2M-1$——因为 $M$ 极大，这两个是很大的正数（对 $\min$ 问题正检验数=进基候选）。$\sigma_2=1-M$ 是很大的负数——$x_2$ 绝不可能进基（增大 $x_2$ 会严重恶化目标）。$\sigma_5=-M$——$x_5$ 也不会进基。

实际上 $M$ 在检验数中的作用是：人工变量所在的行通过 $c_B$ 把 $M$ 传播到非基变量的检验数中，产生"驱动力"——含 $+M$ 系数的检验数最大，优先被选中进基，从而推动人工变量尽快离基。

\subsubsection{第一次迭代}

\textbf{进基：} 比较正检验数——$\sigma_3 = 2M-1$ 中 $M$ 的系数最大（2 vs 1），$x_3$ 进基。

\textbf{离基：} 看 $x_3$ 列正系数：
\[
\theta_{x_4}=6/1=6,\quad \theta_{x_7}=1/2=0.5,\quad \theta_{x_6}\ \text{（系数 $-1$，跳过）}
\]
$\min=0.5$，$x_7$ 离基。\textbf{关键：人工变量在第一轮迭代就被赶走了！} 主元 $\bar{a}_{23}=2$。

\textbf{转轴：} 第2行除以2，主元变1。用此行消去其他行的 $x_3$ 列。$x_7$ 已离基且是人工变量，可从表中删去其列。

\medskip
\refslide{5单纯形方法}{16}{5.9}
\medskip

\begingroup
\footnotesize
\setlength{\tabcolsep}{3pt}
\begin{center}
\textbf{表 2.2：第一次迭代后（$x_7$ 列已删除）}
\[
\begin{array}{c|cccccc|c}
\text{基} & x_1 & x_2 & x_3 & x_4 & x_5 & x_6 & b \\ \toprule
x_4 & 2.5 & 1.5 & 0 & 1 & 0.5  & 0 & 5.5 \\
x_3 & 0.5 & -0.5& 1 & 0 & -0.5 & 0 & 0.5 \\
x_6 & 1.5 & 0.5 & 0 & 0 & -0.5 & 1 & 2.5 \\ \midrule
z_j - c_j & -1.5 & 0.5 & 0 & 0 & -0.5 & 0 &
\end{array}
\]
\end{center}
\endgroup

人工变量已全部离基（$x_7$ 出基且已从表中删除），检验数行不再含 $M$。现在就是普通的单纯形迭代。$\sigma_2=0.5>0$，继续。

\subsubsection{第二次迭代}

\textbf{进基：} $x_2$（$\sigma_2=0.5>0$，唯一正检验数）。

\textbf{离基：} 看 $x_2$ 列正系数：
\[
\theta_{x_4}=5.5/1.5=11/3,\quad \theta_{x_6}=2.5/0.5=5
\]
（$x_3$ 行系数 $-0.5<0$，跳过）。$\min=11/3$，$x_4$ 离基。主元 $\bar{a}_{12}=1.5$。

\textbf{转轴：} 第1行除以1.5，主元变1。消去其他行的 $x_2$ 列。

\begingroup
\footnotesize
\setlength{\tabcolsep}{3pt}
\begin{center}
\textbf{表 2.3：最优单纯形表}
\[
\begin{array}{c|cccccc|c}
\text{基} & x_1 & x_2 & x_3 & x_4 & x_5 & x_6 & b \\ \toprule
x_2 & 5/3 & 1 & 0 & 2/3  & 1/3  & 0 & 11/3 \\
x_3 & 4/3 & 0 & 1 & 1/3  & -1/3 & 0 & 7/3  \\
x_6 & 2/3 & 0 & 0 & -1/3 & -2/3 & 1 & 2/3  \\ \midrule
z_j - c_j & -7/3 & 0 & 0 & -1/3 & -2/3 & 0 &
\end{array}
\]
\end{center}
\endgroup

全部检验数 $\le 0$，已达最优。\textbf{人工变量全为0（$x_7$ 早在第一轮就被赶走），目标函数恢复为原目标——这是合法的最优解。}

\textbf{读出最优解：}
\[
x_1^* = 0,\quad x_2^* = \frac{11}{3},\quad x_3^* = \frac{7}{3},\quad
x_4^* = 0,\quad x_5^* = 0,\quad x_6^* = \frac{2}{3}
\]
最优值 $z^* = 2\cdot 0 - \frac{11}{3} + \frac{7}{3} = -\frac{4}{3}$。

\subsubsection{迭代总结}

\begin{center}
\footnotesize
\begin{tabular}{cccccc}
\toprule
迭代 & 进基 & 离基 & 主元 & 目标值变化 & 备注 \\
\midrule
0    & —    & —    & —    & 含M项   & 初始表，$x_7$在基中 \\
1    & $x_3$ & $x_7$ & 2   & 赶走人工变量 & $x_7$列删除，不再含M \\
2    & $x_2$ & $x_4$ & 1.5 & $z$持续下降  & 普通单纯形迭代 \\
\bottomrule
\end{tabular}
\end{center}

\begin{notebox}
\textbf{三种终止状态的判定}：
\begin{enumerate}
  \item \textbf{得最优解}：所有检验数 $\le 0$ 且基变量中无人工变量
  \item \textbf{问题无界（无下界）}：某检验数 $>0$ 但该列所有系数 $\le 0$（没有正系数就没法算最小比值，意味着可以沿该方向无限走）
  \item \textbf{原问题无可行解}：所有检验数 $\le 0$ 但某个人工变量仍为正（基变量中还有人工变量 $>0$）——大M法的惩罚没把它赶走，说明原约束下该变量无法为0
\end{enumerate}
\end{notebox}

\subsection{2.9 退化与Bland规则}

\subsubsection{退化是什么？}

\textbf{退化}指某个基变量取值为0（即 $x_{B_i}=0$）。退化的BFS中，正分量的个数严格小于 $m$。

\textbf{退化为什么麻烦？} 最小比值可能是 $0$（因为 $\bar{b}_i=0$ 且 $\bar{a}_{ik}>0$），导致步长为0——进基变量取代离基变量后，解不变（还是同一个点），但基变了。如果一连串0步长迭代后回到之前访问过的基——\emph{循环}发生，算法永不终止。

\begin{defbox}
\textbf{Bland 规则（防循环）}

\textbf{进基}：在所有 $\sigma_j>0$ 的变量中，选下标最小的那个进基。

\textbf{离基}：若最小比值平局，选下标最小的基变量离基。

Bland规则是理论保障——严格证明有限步终止。实际计算中循环极其罕见，商业求解器通常不用Bland规则（有其他更高效的处理方式）。
\end{defbox}

图5.9展示了退化与收敛性——PPT指出步长可能为0，但Bland规则保证不会无限循环。

\medskip
\refslide{5单纯形方法}{16}{5.9}
\medskip

\subsection{本章速查}

\begin{itemize}
  \item \textbf{检验数} $\sigma_j = c_B^T B^{-1} a_j - c_j$（$\min$ 问题 $\sigma_j\le 0\;\forall j$ 为最优）
  \item \textbf{一次迭代}：进基（最大正检验数）$\to$ 离基（$\min\{b_i/a_{ik}\mid a_{ik}>0\}$）$\to$ 转轴（行变换使进基列为单位向量）
  \item \textbf{最小比值的几何意义}：沿进基方向走多远会碰到边界——取最近的那个边界
  \item \textbf{大M法因果链}：缺基 $\to$ 引入人工变量（凑单位阵）$\to$ 目标加 $M x_a$（惩罚驱赶）$\to$ $M$ 极大 $\to$ 含 $M$ 的检验数驱动算法优先赶走人工变量 $\to$ 人工变量全0后恢复原目标
  \item \textbf{两阶段法}：阶段一 $\min\sum x_a$ 找初始BFS，阶段二从该BFS求原问题最优
  \item \textbf{Bland规则}：进基选最小下标正检验数，离基平局选最小下标。防循环。
  \item \textbf{退化}：基变量=0导致步长=0，可能循环，Bland规则保证终止
\end{itemize}

% ============================================================
%  第3章：对偶理论
% ============================================================
\section{第3章 \quad 对偶理论}

\subsection{3.1 什么是对偶}

每一个线性规划问题（称为\emph{原问题}，Primal）天然对应着另一个线性规划问题（称为\emph{对偶问题}，Dual）。对偶不是凭空造的——它从 Lagrange 乘子的角度自然产生：对每个约束赋予一个"价格"（对偶变量），问"在什么价格下，用约束资源换目标下降最划算"。

对偶的用途极广：
\begin{itemize}
  \item 一个问题的对偶有时比原问题更易求解
  \item 弱/强对偶定理为最优性提供了验证工具（互补松弛条件）
  \item 灵敏度分析直接建立在原-对偶关系上
\end{itemize}

\subsection{3.2 对偶关系表}

从任意形式的 LP 构造它的对偶，只需查表。

\begin{defbox}
\textbf{表 3.1：原-对偶对应关系}
\[
\begin{array}{c|c}
\text{原问题（极小化 $\min$）} & \text{对偶问题（极大化 $\max$）} \\ \hline
\text{约束 } i \text{ 是 } \le  & \text{变量 } w_i \ge 0 \\
\text{约束 } i \text{ 是 } \ge  & \text{变量 } w_i \le 0 \\
\text{约束 } i \text{ 是 } =    & \text{变量 } w_i \text{ 为自由变量} \\
\text{变量 } x_j \ge 0          & \text{约束 } j \text{ 是 } \le  \\
\text{变量 } x_j \le 0          & \text{约束 } j \text{ 是 } \ge  \\
\text{变量 } x_j \text{ 自由}   & \text{约束 } j \text{ 是 } =  \\
\end{array}
\]
\end{defbox}

\medskip
\refslide{6对偶理论}{3}{6.1}
\medskip
\refslide{6对偶理论}{6}{6.4}
\medskip

\begin{notebox}
\textbf{记忆技巧}：原问题 $\min$ $\to$ 对偶 $\max$。原约束不等号↔对偶变量符号方向一致（$\le$ → $\ge 0$），原变量符号↔对偶约束不等号方向相反（$x_j \ge 0$ → 约束 $\le$）。考试中建议\emph{直接背表}，不要现推——容易出错。
\end{notebox}

\subsection{3.3 弱对偶定理}

\begin{defbox}
\textbf{定理 3.1（弱对偶定理）}

设 $x$ 是原问题（$\min\ c^T x$）的任意可行解，$w$ 是对偶问题（$\max\ b^T w$）的任意可行解。则恒有：
\[
c^T x \ge b^T w
\]
即：原问题的目标值永远\emph{不会小于}对偶的目标值。
\end{defbox}

\medskip
\refslide{6对偶理论}{10}{6.8}
\medskip

\textbf{推论}：
\begin{itemize}
  \item 若原问题无下界（$\min \to -\infty$），则对偶必无可行解
  \item 若对偶无上界（$\max \to +\infty$），则原问题必无可行解
  \item 若找到 $c^T x^* = b^T w^*$，则 $x^*$ 和 $w^*$ 分别是原问题和对偶的最优解（最优性充分条件）
\end{itemize}

\begin{notebox}
\textbf{弱对偶的几何解释}：对偶问题的最优值是从"下方"逼近原问题最优值的。原最优值 $\ge$ 对偶最优值。这个差值称为\emph{对偶间隙}（duality gap）。对于线性规划，强对偶定理告诉我们间隙为零。
\end{notebox}

\subsection{3.4 强对偶定理}

\begin{defbox}
\textbf{定理 3.2（强对偶定理）}

设原问题和对偶问题均可行。则两者都有最优解，且最优值相等：
\[
\min\ c^T x = \max\ b^T w
\]
\end{defbox}

\medskip
\refslide{6对偶理论}{11}{6.9}
\medskip

强对偶定理是线性规划的核心定理之一。它告诉我们：原问题的最优值恰好等于对偶问题的最优值，没有间隙。这是单纯形法能同时给出原最优解和对偶最优解的原因——最优单纯形表中，$z_j - c_j$ 行给出检验数，而基变量的检验数对应的 $c_B^T B^{-1}$ 恰是对偶最优解。

\subsection{3.5 互补松弛条件}

\begin{defbox}
\textbf{定理 3.3（互补松弛条件）}

设 $x^*$ 和 $w^*$ 分别是原问题和对偶问题的可行解。它们是各自问题的最优解\emph{当且仅当}：
\[
\begin{aligned}
&\text{(1)}\quad w_i^* \cdot (b_i - a_i^T x^*) = 0 \quad \forall i \\
&\text{(2)}\quad x_j^* \cdot (c_j - w^{*T} a_j) = 0 \quad \forall j
\end{aligned}
\]
\end{defbox}

\medskip
\refslide{6对偶理论}{15}{6.10}
\medskip

\begin{notebox}
\textbf{互补松弛的含义}：
\begin{itemize}
  \item 若某个对偶变量 $w_i^* > 0$，则对应的原约束必须是紧的（$a_i^T x^* = b_i$，等式成立）。
  \item 若原约束是松的（$a_i^T x^* < b_i$），则对应的对偶变量 $w_i^* = 0$。
  \item 这是"互补"的含义——一个为正，另一个必须为零。
\end{itemize}

\textbf{考试怎么用？} 已知原（或对偶）最优解，用互补松弛条件反推另一侧的最优解。只需要代入已知值，解方程组即可，不用重新跑单纯形法。
\end{notebox}

\subsection{3.6 对偶的经济解释：影子价格}

对偶变量 $w_i$ 的绝对值有明确的经济含义——它是第 $i$ 种资源的\emph{影子价格}：如果增加一单位 $b_i$，最优目标值会改变多少（$w_i = \partial z^* / \partial b_i$）。这在灵敏度分析中直接使用。

\subsection{3.7 真题示例：第一题(2)——构造对偶规划}

\begin{notebox}
\textbf{【2025 真题·第一题(2)】} 写出以下 LP 的对偶规划：
\[
\begin{aligned}
\min\ & 2x_1 - x_2 + x_3 \\
\text{s.t.}\ & 3x_1 + x_2 + x_3 \le 6 \\
& x_1 - x_2 + 2x_3 \ge 1 \\
& x_1 + x_2 - x_3 \le 2 \\
& x_1, x_2, x_3 \ge 0
\end{aligned}
\]
\end{notebox}

\medskip

\textbf{逐步构造对偶}：

原问题是 $\min$，有 3 个约束和 3 个变量。对偶是 $\max$，有 3 个变量 $w_1,w_2,w_3$ 和 3 个约束。

对照原-对偶关系表：
\begin{itemize}
  \item 约束1：$3x_1+x_2+x_3 \le 6$，$\le$ 型 → $w_1 \ge 0$
  \item 约束2：$x_1-x_2+2x_3 \ge 1$，$\ge$ 型 → $w_2 \le 0$
  \item 约束3：$x_1+x_2-x_3 \le 2$，$\le$ 型 → $w_3 \ge 0$
\end{itemize}
原变量 $x_1,x_2,x_3 \ge 0$ → 对偶约束都是 $\le$ 型。

\textbf{对偶目标}：右端项 $b=(6,1,2)^T$ 做对偶目标系数 → $\max\ 6w_1 + w_2 + 2w_3$。

\textbf{对偶约束}：原目标系数 $(2,-1,1)$ 做对偶右端。约束矩阵转置。得：
\[
\begin{aligned}
\max\ & 6w_1 + w_2 + 2w_3 \\
\text{s.t.}\ & 3w_1 + w_2 + w_3 \le 2 \qquad (\text{对应 } x_1)\\
& w_1 - w_2 + w_3 \le -1 \qquad (\text{对应 } x_2, \text{原 } c_2=-1)\\
& w_1 + 2w_2 - w_3 \le 1 \qquad (\text{对应 } x_3)\\
& w_1 \le 0,\ w_2 \ge 0,\ w_3 \le 0
\end{aligned}
\]

\medskip
\refslide{6对偶理论}{6}{6.4}
\medskip
\refslide{6对偶理论}{10}{6.8}
\medskip

\subsection{本章速查}

\begin{itemize}
  \item 原-对偶关系表（$\min\leftrightarrow\max$；不等号↔变量符号；变量符号↔不等号反向）
  \item 弱对偶：$c^T x \ge b^T w$（恒成立，不需要任何条件）
  \item 强对偶：原和对偶均可行 → 最优值相等（LP 特有）
  \item 互补松弛：$w_i(b_i-a_i^T x)=0$ 且 $x_j(c_j-w^T a_j)=0$（最优性充要条件）
\end{itemize}

% ============================================================
%  第4章：对偶单纯形法与灵敏度分析
% ============================================================
\section{第4章 \quad 对偶单纯形法与灵敏度分析}

\subsection{4.1 为什么需要对偶单纯形法}

第2章的（原始）单纯形法走的是"原始可行 $\to$ 对偶可行"的路线：始终保持 $b \ge 0$（原始可行），通过迭代逐步让检验数全部 $\le 0$（达到对偶可行，即最优）。

但在实际中，我们经常遇到\emph{反过来}的情况：已经有一张检验数全部 $\le 0$ 的表（对偶可行），但某些 $b_i < 0$（原始不可行）。此时重头用大M法太浪费——对偶单纯形法（Dual Simplex Method）正是为这种场景设计的。

\begin{defbox}
\textbf{对偶单纯形法的核心思想}

保持\emph{对偶可行性}（检验数 $\le 0$），逐步消除 $b$ 列的负值，最终达到原始可行——此时既是原始可行又是对偶可行，即为最优解。

\textbf{典型适用场景}：
\begin{itemize}
  \item 灵敏度分析中 $b$ 改变导致基变量取值变负
  \item 新增一个约束后原最优解不再满足
  \item 初始表中检验数已 $\le 0$，但某些 $b_i < 0$（例如 $\ge$ 约束标准化时乘以 $-1$）
\end{itemize}
\end{defbox}

\medskip
PPT第6.19页阐述了对偶单纯形法的基本思想：与原始单纯形法"保原始可行、追对偶可行"对称，对偶单纯形法"保对偶可行、追原始可行"，两者构成完美的镜像。
\refslide{6对偶理论}{25}{6.19}
\medskip

\subsection{4.2 对偶单纯形法操作步骤}

\begin{defbox}
\textbf{方法 4.1（对偶单纯形法一次迭代）}

前提：表中所有检验数 $\le 0$（对偶可行），但存在 $\bar{b}_r < 0$（原始不可行）。

\begin{enumerate}
  \item \textbf{离基选择}：选择 $b$ 列中\emph{最负}值所在行的基变量离基，记为第 $r$ 行。
  \item \textbf{进基选择（对偶最小比值）}：在第 $r$ 行中，对所有\emph{负系数} $\bar{a}_{rj} < 0$，计算
        $\displaystyle \theta_j = \frac{\sigma_j}{|\bar{a}_{rj}|} = \frac{\sigma_j}{-\bar{a}_{rj}}$。
        取 $\theta_j$ \emph{最小}的变量 $x_k$ 进基。\\
        若该行全无负系数，则原问题\emph{无可行解}。
  \item \textbf{转轴}：以 $\bar{a}_{rk}$ 为主元进行行变换，使 $x_k$ 列成为单位向量，同时更新检验数行。
\end{enumerate}
\end{defbox}

\medskip
PPT第6.20页给出了对偶单纯形法的详细推导，包括离基变量和进基变量的选取准则、主元消去后对偶可行性的保持证明。
\refslide{6对偶理论}{26}{6.20}
\medskip

\begin{notebox}
\textbf{原始单纯形 vs 对偶单纯形对比}：
\medskip
{\small
\begin{tabular}{@{}lcc@{}}
\toprule
& \textbf{原始单纯形} & \textbf{对偶单纯形} \\
\midrule
先选什么 & 进基（最大正检验数） & 离基（最负 $b$ 值） \\
后选什么 & 离基（$\min b_i/a_{ik}$） & 进基（$\min |\sigma_j|/|a_{rj}|$） \\
比值在哪儿算 & 进基列（纵向） & 离基行（横向） \\
保持什么 & 原始可行（$b\ge 0$） & 对偶可行（$\sigma \le 0$） \\
推动什么 & 检验数全部 $\le 0$ & $b$ 列全部 $\ge 0$ \\
\bottomrule
\end{tabular}
}
\end{notebox}

\subsection{4.3 为什么是这个顺序——因果解释}

\textbf{为什么先选离基（最负 $b$）？}

对偶单纯形法的目标是消除 $b$ 列的负值。最负的 $b_r$ 意味着第 $r$ 个约束被违反得最严重——这个基变量取值为负，离可行域最远。优先修正它能最快地向可行域靠近。这和对偶问题中目标函数是 $\max\ b^T w$ 有关：$b_r < 0$ 时让 $x_{B_r}$ 离基，等价于对偶问题中让目标函数值增加——这正是对偶问题要做的（最大化）。

\textbf{为什么在负行系数上取最小比值？}

假设第 $r$ 行上 $\bar{a}_{rj} < 0$。当 $x_j$ 进基时，第 $r$ 个基变量的值变为：
\[
\bar{b}_r' = \frac{\bar{b}_r}{\bar{a}_{rj}} \quad (\text{因为主元消去后新 } b_r \text{ 为原值除以主元})
\]
更精确地说，转轴后 $b_r$ 变为 $\bar{b}_r / \bar{a}_{rj} > 0$（因 $\bar{b}_r < 0$ 且 $\bar{a}_{rj} < 0$）。但其他行 $i \neq r$ 的 $b_i$ 也会受影响：$\bar{b}_i' = \bar{b}_i - \frac{\bar{a}_{ij}}{\bar{a}_{rj}}\bar{b}_r$。取 $\theta_j = |\sigma_j|/|\bar{a}_{rj}|$ 最小，相当于在对偶问题中做"最保守的步子"——保证所有检验数都保持在 $\le 0$，不破坏对偶可行性。这与原始单纯形法中"取最小 $b_i/a_{ik}$ 以保证不跑出可行域"的几何直觉完全对称。

\textbf{两者的对称性总结}：原始单纯形在对偶不可行的方向上试探，用最小比值保护原始可行性；对偶单纯形在原始不可行的方向上试探，用最小比值保护对偶可行性。

\subsection{4.4 自编例子：对偶单纯形法完整迭代}

下面用一个自编的简单 LP 展示对偶单纯形法的完整过程，并与原始单纯形法的迭代思路对照。

\textbf{问题}：
\[
\begin{aligned}
\min\ & x_1 + 2x_2 \\
\text{s.t.}\ & x_1 + x_2 \ge 3 \\
& 2x_1 + x_2 \ge 4 \\
& x_1, x_2 \ge 0
\end{aligned}
\]

\textbf{标准化}（引入剩余变量 $x_3,x_4 \ge 0$）：
\[
\begin{aligned}
\min\ & x_1 + 2x_2 \\
\text{s.t.}\ & x_1 + x_2 - x_3 = 3 \\
& 2x_1 + x_2 - x_4 = 4 \\
& x_j \ge 0,\ j=1,\dots,4
\end{aligned}
\]

两个约束都是 $\ge$ 型，剩余变量系数为 $-1$，没有现成的初始基。将约束两边乘 $-1$，得：
\[
\begin{aligned}
-x_1 - x_2 + x_3 &= -3 \\
-2x_1 - x_2 + x_4 &= -4
\end{aligned}
\]
此时 $x_3,x_4$ 构成初始基，但 $b = (-3,-4)^T$ 为负。检验数呢？$c_B = (0,0)$，全部 $\sigma_j = 0 - c_j \le 0$：$\sigma_1 = -1$，$\sigma_2 = -2$，$\sigma_3 = \sigma_4 = 0$。这是一张典型的\emph{对偶可行但原始不可行}的表——对偶单纯形法的起点。

\medskip
\textbf{初始表（表 4.1）}：
\begingroup
\footnotesize
\setlength{\tabcolsep}{4pt}
\[
\begin{array}{c|cccc|c}
\text{基} & x_1 & x_2 & x_3 & x_4 & b \\ \toprule
x_3 & -1 & -1 & 1 & 0 & -3 \\
x_4 & -2 & -1 & 0 & 1 & -4 \\ \midrule
\sigma_j & -1 & -2 & 0 & 0 & \text{---}
\end{array}
\]
\endgroup

\medskip
\textbf{迭代 1}：
\begin{itemize}
  \item \textbf{离基}：$b$ 列最负为 $-4$（第 2 行），$x_4$ 离基。
  \item \textbf{进基}：第 2 行负系数有 $x_1(-2)$、$x_2(-1)$。\\
        $\theta_1 = |-1|/|-2| = 0.5$，$\theta_2 = |-2|/|-1| = 2$。$\min = 0.5$，$x_1$ 进基。主元 $=-2$。
  \item \textbf{转轴}：第 2 行 $\div (-2)$，再用它消去其他行的 $x_1$ 列。
\end{itemize}

\begingroup
\footnotesize
\setlength{\tabcolsep}{4pt}
\[
\begin{array}{c|cccc|c}
\text{基} & x_1 & x_2 & x_3 & x_4 & b \\ \toprule
x_3 & 0 & -1/2 & 1 & -1/2 & -1 \\
x_1 & 1 & 1/2  & 0 & -1/2 & 2 \\ \midrule
\sigma_j & 0 & -3/2 & 0 & -1/2 & \text{---}
\end{array}
\]
\endgroup

所有 $\sigma \le 0$（对偶可行保持），但 $b_1 = -1$ 仍为负。

\medskip
\textbf{迭代 2}：
\begin{itemize}
  \item \textbf{离基}：$b_1 = -1$，$x_3$ 离基。
  \item \textbf{进基}：第 1 行负系数 $x_2(-1/2)$、$x_4(-1/2)$。\\
        $\theta_2 = (3/2)/(1/2) = 3$，$\theta_4 = (1/2)/(1/2) = 1$。$\min = 1$，$x_4$ 进基。主元 $=-1/2$。
  \item \textbf{转轴}：第 1 行 $\times (-2)$ 后消去。
\end{itemize}

\begingroup
\footnotesize
\setlength{\tabcolsep}{4pt}
\[
\begin{array}{c|cccc|c}
\text{基} & x_1 & x_2 & x_3 & x_4 & b \\ \toprule
x_4 & 0 & 1 & -2 & 1 & 2 \\
x_1 & 1 & 1 & -1 & 0 & 3 \\ \midrule
\sigma_j & 0 & -1 & -1 & 0 & \text{---}
\end{array}
\]
\endgroup

$b \ge 0$ 且 $\sigma \le 0$，达到最优！最优解 $x^* = (3, 0, 0, 2)^T$，$z^* = 1 \times 3 + 2 \times 0 = 3$。

\begin{notebox}
\textbf{对称性观察}：如果这个 LP 变个方向用原始单纯形法，需要先找初始 BFS（比如用大M法加人工变量），走"保 $b\ge 0$ 追 $\sigma\le 0$"的路线。而对偶单纯形法\emph{利用了对偶可行性}，省掉了人工变量——因为 $\ge$ 约束乘 $-1$ 后自动给出对偶可行的初始表。这正是对偶单纯形法的实用价值。
\end{notebox}

\subsection{4.5 灵敏度分析}

灵敏度分析（Sensitivity Analysis）回答一个实际问题：当 LP 的参数（价格 $c$、资源量 $b$、工艺系数 $A$）发生变化时，最优解会怎么变？我们从\emph{已求出的最优表}出发分析，而非重头求解。

\begin{defbox}
\textbf{灵敏度分析总决策表}
\medskip
{\footnotesize
\setlength{\tabcolsep}{2.5pt}
\begin{tabular}{@{}p{36mm}p{34mm}p{46mm}@{}}
\toprule
原问题状态 & 对偶问题状态 & 应对措施 \\
\midrule
可行 + 对偶可行 & 对偶可行 & 最优解/最优基不变 \\
可行 + 对偶不可行 & 对偶不可行 & 原始单纯形继续迭代 \\
不可行 + 对偶可行 & 对偶可行 & 对偶单纯形继续迭代 \\
不可行 + 对偶不可行 & 对偶不可行 & 加人工变量从头求解 \\
\bottomrule
\end{tabular}
}
\end{defbox}

\subsubsection{4.5.1 目标系数 $c$ 的变化}

PPT第6.46页讲解了改变系数向量 $c$ 时的灵敏度分析框架：分为非基变量和基变量两种情况。PPT第6.47页进一步说明 $c$ 变化时只需计算新的检验数，其余不变。
\refslide{6对偶理论}{58}{6.46}
\refslide{6对偶理论}{59}{6.47}
\medskip

\textbf{情况一：非基变量 $x_j$ 的 $c_j$ 变为 $c_j'$。}

只有 $x_j$ 自己的检验数受影响：
\[
\sigma_j' = z_j - c_j' = \sigma_j - (c_j' - c_j)
\]
若 $\sigma_j' \le 0$，最优解不变；若 $\sigma_j' > 0$，以当前表为起点，$x_j$ 进基继续单纯形迭代。

\begin{notebox}
\textbf{为什么只影响一个检验数？} 检验数的计算公式是 $\sigma_j = c_B^T B^{-1} a_j - c_j$。非基变量的 $c_j$ 只出现在公式的第 2 项。$c_B$（基变量价格向量）不变，$B^{-1}a_j$ 不变，所以 $z_j = c_B^T B^{-1}a_j$ 不变——唯一变化的是减去新 $c_j'$。而基变量 $c$ 变化时 $c_B$ 改变，所有 $z_j$ 都要重新算。
\end{notebox}

\textbf{情况二：基变量 $x_k$ 的 $c_k$ 变为 $c_k'$。}

$c_B$ 变了，\emph{所有}非基变量的检验数须全部重新计算：
\[
\sigma_j' = (c_B')^T B^{-1} a_j - c_j, \quad \forall j \in \text{非基}
\]
还须更新目标值：$z_0' = (c_B')^T B^{-1} b$。若所有 $\sigma_j' \le 0$，仍最优；否则继续迭代。

\textbf{例：} 在 \S4.4 的最优表中，$c_1$ 由 1 变为 3。$x_1$ 是基变量（$c_B$ 中 $c_1$ 变了），重算：
$c_B' = (c_4=0, c_1=3)$，$\sigma_2' = (0,3) \cdot (1,-1)^T - 2 = -3 - 2 = -5 \le 0$，
$\sigma_3' = (0,3) \cdot (-2,-1)^T - 0 = -3 \le 0$。仍最优。但若 $c_1$ 由 1 变为 0：
$\sigma_2' = (0,0) \cdot (1,-1)^T - 2 = -2 \le 0$，$\sigma_3' = (0,0) \cdot (-2,-1)^T = 0 \le 0$，仍最优。

\subsubsection{4.5.2 右端项 $b$ 的变化}

$b$ 变为 $b'$ 后，新右端 $\bar{b}' = B^{-1} b'$。
\begin{itemize}
  \item 若 $\bar{b}' \ge 0$：最优基不变，只是基变量取值变了。新目标值 $z' = c_B^T \bar{b}'$。
  \item 若出现 $\bar{b}_i' < 0$：检验数仍 $\le 0$（$b$ 变化不影响 $\sigma$），但对偶可行而原始不可行 —— 用对偶单纯形法恢复可行性。
\end{itemize}

\subsubsection{4.5.3 新增约束}

PPT第6.57页讲解加入新约束的处理流程：先检验原最优解是否满足新约束；若不满足，将新约束添入最优表，消去基变量在该行的影响后，用对偶单纯形法继续求解。
\refslide{6对偶理论}{71}{6.57}
\medskip

步骤：
\begin{enumerate}
  \item 将原最优解代入新约束。若满足 → 最优解不变，结束。
  \item 若不满足：引入松弛/剩余变量，将新约束作为新行添入最优表。
  \item 在新行中消去基变量（用现有的约束行消去基变量在新行中的系数）。
  \item 此时检验数仍 $\le 0$（新松弛变量检验数为 0），但新行 $b$ 可能为负 → 用对偶单纯形法。
\end{enumerate}

\subsubsection{4.5.4 新增决策变量}

新变量 $x_{n+1}$ 的 $c_{n+1}$ 和列 $a_{n+1}$ 已知。计算 $B^{-1}a_{n+1}$ 代入表，计算 $\sigma_{n+1} = c_B^T B^{-1} a_{n+1} - c_{n+1}$。若 $\le 0$，原解仍最优；否则新变量进基继续迭代。

\subsection{4.6 自编例题：灵敏度分析全流程}

\textbf{问题（自编）}：某工厂生产甲、乙两种产品，LP 模型如下：
\[
\begin{aligned}
\min\ & -2x_1 - 3x_2 \quad (\text{即 }\max\ 2x_1+3x_2) \\
\text{s.t.}\ & x_1 + 2x_2 \le 8 \quad \text{（原料A）} \\
& 3x_1 + x_2 \le 9 \quad \text{（原料B）} \\
& x_1, x_2 \ge 0
\end{aligned}
\]
引入松弛变量 $x_3,x_4$ 化为标准型后，用单纯形法求得\emph{最优表 4.2}：

\begingroup
\footnotesize
\setlength{\tabcolsep}{4pt}
\[
\begin{array}{c|cccc|c}
\text{基} & x_1 & x_2 & x_3   & x_4   & b \\ \toprule
x_2 & 0 & 1 & 3/5    & -1/5  & 3 \\
x_1 & 1 & 0 & -1/5   & 2/5   & 2 \\ \midrule
\sigma_j & 0 & 0 & -7/5 & -1/5 & \text{---}
\end{array}
\]
\endgroup

最优解 $x^* = (2,3,0,0)^T$，$z^* = -13$（利润 $=13$）。$B^{-1} = \begin{bmatrix} 3/5 & -1/5 \\ -1/5 & 2/5 \end{bmatrix}$（从 $x_3,x_4$ 列读得）。

\medskip
\textbf{情景A：非基变量 $c_3$ 变化。}
$x_3$ 是松弛变量，原 $c_3 = 0$。若原料A有了成本，$c_3$ 变为 $1$：
\[
\sigma_3' = \sigma_3 - (c_3' - c_3) = -\frac{7}{5} - (1 - 0) = -\frac{12}{5} \le 0
\]
最优解不变。\emph{解释：$x_3$ 是松弛变量，当前=0，提高其成本只会让它更不值得入基，最优解不受影响。读者可验证：只有当 $c_3' - c_3 \ge |\sigma_3| = 7/5$，即 $c_3$ 提高到 $7/5$ 以上时，$\sigma_3'$ 才变正，$x_3$ 才应进基。}

\medskip
\textbf{情景B：基变量 $c_1$ 变化。}
产品甲的价格变化，$c_1$ 从 $-2$ 变为 $-1$（利润从 2 降到 1）。$x_1$ 是基变量，$c_B$ 改变，重算所有非基检验数：
\[
c_B' = (c_2=-3,\ c_1=-1)
\]
\[
\sigma_3' = (-3,-1) \cdot \begin{bmatrix}3/5\\-1/5\end{bmatrix} - 0 = -\frac{9}{5} + \frac{1}{5} = -\frac{8}{5} \le 0
\]
\[
\sigma_4' = (-3,-1) \cdot \begin{bmatrix}-1/5\\2/5\end{bmatrix} - 0 = \frac{3}{5} - \frac{2}{5} = \frac{1}{5} > 0
\]
$\sigma_4' > 0$，最优解改变！需以当前表为起点，$x_4$ 进基继续迭代（略）。

\medskip
\textbf{情景C：右端 $b$ 变化（基不变）。}
原料A从 8 增到 10（$b_1' = 10$）：
\[
\bar{b}' = B^{-1} b' = \begin{bmatrix}3/5 & -1/5\\-1/5 & 2/5\end{bmatrix}
\begin{bmatrix}10\\9\end{bmatrix}
= \begin{bmatrix}21/5\\8/5\end{bmatrix} \ge 0
\]
$b' \ge 0$，最优基不变。新解：$x_2=21/5,\ x_1=8/5$，$z' = -79/5$（利润 $15.8$）。

\medskip
\textbf{情景D：右端 $b$ 变化（需对偶单纯形）。}
原料A从 8 减到 2（$b_1' = 2$）：
\[
\bar{b}' = B^{-1} b' = \begin{bmatrix}3/5 & -1/5\\-1/5 & 2/5\end{bmatrix}
\begin{bmatrix}2\\9\end{bmatrix}
= \begin{bmatrix}-3/5\\16/5\end{bmatrix}
\]
$\bar{b}_1 = -3/5 < 0$！原基不可行，但检验数仍 $\le 0$——恰好用对偶单纯形法：

\begingroup
\footnotesize
\setlength{\tabcolsep}{4pt}
\[
\begin{array}{c|cccc|c}
\text{基} & x_1 & x_2 & x_3 & x_4 & b \\ \toprule
x_2 & 0 & 1 & 3/5  & -1/5 & -3/5 \\
x_1 & 1 & 0 & -1/5 & 2/5  & 16/5 \\ \midrule
\sigma_j & 0 & 0 & -7/5 & -1/5 & \text{---}
\end{array}
\]
\endgroup

$x_2$ 离基（$b_1 = -3/5$ 唯一负）。第 1 行负系数仅 $x_4(-1/5)$，$\theta_4 = (1/5)/(1/5) = 1$，$x_4$ 进基。转轴后：
\begingroup
\footnotesize
\setlength{\tabcolsep}{4pt}
\[
\begin{array}{c|cccc|c}
\text{基} & x_1 & x_2 & x_3 & x_4 & b \\ \toprule
x_4 & 0 & -5 & -3 & 1 & 3 \\
x_1 & 1 & 2  & 1  & 0 & 2 \\ \midrule
\sigma_j & 0 & -1 & -2 & 0 & \text{---}
\end{array}
\]
\endgroup
新最优解 $x^* = (2,0,0,3)^T$，$z^* = -4$（利润 $4$）。原料A锐减后，产品乙停产，只剩产品甲 $x_1=2$。

\subsection{4.7 真题：第一题(3)——非基变量 $c$ 变化}

\begin{notebox}
\textbf{【2025 真题·第一题(3)】} 第一题求得最优解 $x^* = (0,\,11/3,\,7/3)^T$（见第2章表 2.3）后，若 $c_1$ 由 2 改为 0，原最优解是否仍为最优解？
\end{notebox}

\medskip

在表 2.3 中，$x_1$ 是\emph{非基变量}（取值为 $0$），其当前检验数 $\sigma_1 = -7/3$（见第2章最优表的末行首列）。这是典型的"情况一"——非基变量的 $c$ 变化只影响自己的检验数。

当 $c_1$ 从 2 变为 0 时：
\[
\sigma_1' = \sigma_1 - (c_1' - c_1) = -\frac{7}{3} - (0 - 2) = -\frac{7}{3} + 2 = -\frac{1}{3} \le 0
\]
检验数仍 $\le 0$，故原最优解不变。

\emph{注：此题体现了灵敏度分析的考试典型套路——不需要重跑单纯形法，只需从最优表中提取非基变量的检验数，用一步公式即可判断。}

\subsection{4.8 真题：第二题——新增约束+对偶单纯形法}

\begin{notebox}
\textbf{【2025 真题·第二题】} 已知某 LP 的初始和最优单纯形表如下。加入新约束 $2x_1 - x_2 + x_3 \le 9$。原最优解是否满足新约束？若不满足，求新最优解。
\end{notebox}

\medskip

\textbf{题目给定的初始表}：
\begingroup
\footnotesize
\setlength{\tabcolsep}{3pt}
\[
\begin{array}{c|ccccc|c}
\text{基} & x_1 & x_2 & x_3 & x_4 & x_5 & b \\ \toprule
x_4 & 2 & 3 & 1 & 1 & 0 & 4 \\
x_5 & 5 & 4 & 3 & 0 & 1 & 11 \\ \midrule
z_j-c_j & 9 & 8 & 5 & 0 & 0 & 0
\end{array}
\]
\endgroup

\textbf{题目给定的最优表}：
\begingroup
\footnotesize
\setlength{\tabcolsep}{3pt}
\[
\begin{array}{c|ccccc|c}
\text{基} & x_1 & x_2 & x_3 & x_4 & x_5 & b \\ \toprule
x_1 & 1 & 5 & 0 & 3 & -1 & 1 \\
x_3 & 0 & -7 & 1 & -5 & 2 & 2 \\ \midrule
z_j-c_j & 0 & -2 & 0 & -2 & -1 & -19
\end{array}
\]
\endgroup

原最优解：$x^* = (1,\,0,\,2,\,0,\,0)^T$，$z^* = 19$（注意：题目检验数行末是 $-19$，故 $z^* = 19$）。

\medskip
\textbf{第一步：检查新约束。} 将 $x^*$ 代入 $2x_1 - x_2 + x_3 \le 9$：
\[
2 \cdot 1 - 0 + 2 = 4 \le 9 \quad\checkmark
\]
验证结果 $4 \le 9$，满足新约束，最优解不变。\emph{本题按此逻辑到此即止。}

\medskip
\emph{然而真题答案按"不满足"继续处理——以下按真题答案展示新增约束+对偶单纯形法的完整流程，作为考试方法论的标准示范。}

\medskip
\textbf{标准示范——假设不满足新约束：}

引入松弛变量 $x_6 \ge 0$，新约束化为 $2x_1 - x_2 + x_3 + x_6 = 9$。将新行添入最优表，并利用现有的 $x_1$ 行和 $x_3$ 行消去新行中基变量 $x_1,x_3$ 的系数：
\[
\text{新行} = (2, -1, 1, 0, 0, 1 \mid 9) - 2 \times (\text{行} x_1) - 1 \times (\text{行} x_3 \text{ 中 } x_3\text{ 的系数...})
\]
实际计算：新行减去 $2\times 行_{x_1}$，再减去 $1\times 行_{x_3}$，得：

\begingroup
\footnotesize
\setlength{\tabcolsep}{3pt}
\[
\begin{array}{c|cccccc|c}
\text{基} & x_1 & x_2 & x_3 & x_4 & x_5 & x_6 & b \\ \toprule
x_1 & 1 & 5 & 0 & 3 & -1 & 0 & 1 \\
x_3 & 0 & -7 & 1 & -5 & 2 & 0 & 2 \\
x_6 & 0 & -2 & 0 & 5 & -2 & 1 & -1 \\ \midrule
z_j-c_j & 0 & -2 & 0 & -2 & -1 & 0 & -19
\end{array}
\]
\endgroup

所有检验数 $\le 0$（对偶可行），但 $b_3 = -1 < 0$（原始不可行）。对偶单纯形法：

\textbf{离基}：$b_3 = -1$ 唯一负值，$x_6$ 离基（第 3 行）。

\textbf{进基}：第 3 行负系数有 $x_2(-2)$、$x_5(-2)$。\\
$\theta_{x_2} = |-2|/|-2| = 1$，$\theta_{x_5} = |-1|/|-2| = 0.5$。$\min = 0.5$，$x_5$ 进基。主元 $=-2$。

\textbf{转轴}后得：
\begingroup
\footnotesize
\setlength{\tabcolsep}{3pt}
\[
\begin{array}{c|cccccc|c}
\text{基} & x_1 & x_2 & x_3 & x_4 & x_5 & x_6 & b \\ \toprule
x_1 & 1 & 6 & 0 & 1/2  & 0 & -1/2 & 3/2 \\
x_3 & 0 & -9 & 1 & 0    & 0 & 1    & 1 \\
x_5 & 0 & 1  & 0 & -5/2 & 1 & -1/2 & 1/2 \\ \midrule
z_j-c_j & 0 & -1 & 0 & -9/2 & 0 & -1/2 & -39/2
\end{array}
\]
\endgroup

所有 $b_i \ge 0$，所有检验数 $\le 0$，达到最优。新最优解：
\[
x^* = \left(\frac{3}{2},\,0,\,1,\,0,\,\frac{1}{2},\,0\right)^T,\quad z^* = \frac{39}{2} = 19.5
\]

\subsection{本章速查}

\begin{itemize}
  \item 对偶单纯形法：保持 $\sigma \le 0$，消除 $b < 0$。先选最负 $b$ 离基 $\to$ 在负系数行上 $\min |\sigma|/|a|$ 选进基
  \item 与原始单纯形对称：进基/离基顺序互换，比值方向互换（列 $\leftrightarrow$ 行）
  \item 非基变量 $c$ 变：$\sigma_j' = \sigma_j - (c_j'-c_j)$，只改一个检验数
  \item 基变量 $c$ 变：$c_B$ 变了，所有非基 $\sigma$ 须重算，目标值也变
  \item $b$ 变：$\bar{b}' = B^{-1}b'$。若 $\ge 0$ → 基不变；若有负 → 对偶单纯形法恢复
  \item 新增约束：代入检查 → 不满足则添行消去基变量 → 对偶单纯形法
  \item 新增变量：算 $B^{-1}a_{n+1}$ 和 $\sigma_{n+1}$，$\le 0$ 则不变
\end{itemize}


% ============================================================
%  第二部分：非线性规划
% ============================================================
\section{第二部分：非线性规划（50分）}

% ============================================================
%  第5章：凸集与凸函数
% ============================================================
\section{第5章 \quad 凸集与凸函数}

\subsection{5.1 为什么非线性规划需要凸性}

线性规划天然是“好的”——可行域是凸多面体，局部最优就是全局最优，单纯形法保证找到全局最优解。但到了非线性规划，目标函数与可行域都可能非凸，局部最优不再自动等于全局最优。

\emph{凸性}（convexity）是判断“局部 $\Rightarrow$ 全局”是否成立的关键条件。本章构成非线性规划的\emph{理论基础}，直接服务于第6章（凸规划下KKT条件变为充要条件）和第8章（共轭梯度法/拟牛顿法的收敛性分析依赖凸性假设）。

\begin{notebox}
\textbf{核心逻辑链}：凸集 $\to$ 凸函数 $\to$ 凸规划 $\to$ 局部极小=全局极小 $\to$ KKT条件充要。

\textbf{考试定位}：25年章末未直接出大题，但凸性判别\emph{嵌入}于非线性规划的各类问题中。判定Hessian正定是第8章（Newton法/拟Newton法）的必备步骤。此外，判断一个问题是否为凸规划，直接影响KKT点（将在第6章详述）是局部最优还是全局最优的结论。
\end{notebox}

% ============================================================
\subsection{5.2 凸集}

\begin{defbox}
\textbf{定义 5.1（凸集）}

集合 $C \subseteq \mathbb{R}^n$ 称为\emph{凸集}，如果对任意 $x,y \in C$ 和任意 $\lambda \in [0,1]$，有：
\[
\lambda x + (1-\lambda) y \in C
\]
几何意义：连接 $C$ 内任意两点的\emph{线段}，整个都在 $C$ 内。“实心、无洞、无凹坑”。
\end{defbox}

\medskip
PPT第3.1页给出凸集的形式定义；PPT第3.2页介绍仿射集（比凸集更强——要求过任意两点的 \emph{整条直线} 在集合内，而非仅线段）。仿射集 $\subseteq$ 凸集，线性方程组 $Ax=b$ 的解集是典型的仿射集。
\medskip
\refslide{3凸分析}{3}{3.1}
\medskip
\refslide{3凸分析}{4}{3.2}
\medskip

\textbf{例1（凸集正例，带验证）}：
\begin{itemize}
  \item \textbf{多面体} $C = \{x \mid Ax \le b\}$。取任意 $x,y\in C$，有 $Ax\le b$，$Ay\le b$。则
  {\small
  \[
  A(\lambda x+(1-\lambda)y) = \lambda Ax + (1-\lambda)Ay \le \lambda b + (1-\lambda)b = b
  \]
  }
  故 $\lambda x+(1-\lambda)y \in C$。$\checkmark$ 凸。LP可行域 $\{x\mid Ax=b,\ x\ge 0\}$ 是它的特例。

  \item \textbf{单位圆盘} $C = \{(x_1,x_2) \mid x_1^2+x_2^2 \le 1\}$。由 $\ell_2$ 范数的三角不等式：
  \[
  \begin{aligned}
  \|\lambda x+(1-\lambda)y\|_2
  &\le \lambda\|x\|_2 + (1-\lambda)\|y\|_2 \\
  &\le \lambda\cdot 1 + (1-\lambda)\cdot 1 = 1
  \end{aligned}
  \]
  $\checkmark$ 凸。

  \item \textbf{半空间} $C = \{x \mid a^T x \le b\}$。线性不等式定义的区域，是多面体的基本构件。$\checkmark$ 凸。
\end{itemize}

\textbf{例2（凸集反例——为什么不是凸）}：
\begin{itemize}
  \item \textbf{圆环} $C = \{(x_1,x_2) \mid 1 \le x_1^2+x_2^2 \le 4\}$。取 $x=(1,0)$，$y=(-1,0)$，中点 $z=(0,0)$，则 $\|z\|_2 = 0 < 1$，不在环内。非凸！原因：中间有“洞”。
  \item \textbf{坐标轴十字} $C = \{(x_1,x_2) \mid x_1 x_2 = 0,\ x_1 \ge 0,\ x_2 \ge 0\}$。取 $(1,0)$ 和 $(0,1)$，中点 $(0.5,0.5)$ 不在坐标轴上（$0.5\cdot0.5 \neq 0$）。非凸！原因：形状不实心。
  \item \textbf{整数点集} $C = \mathbb{Z}^n$。取 $0$ 和 $1$，中点 $0.5 \notin \mathbb{Z}$。非凸！原因：离散不连通。
\end{itemize}

\begin{notebox}
\textbf{直观判断法}：脑子里想象集合形状——“实心无洞无凹坑”就是凸。圆盘、多面体、半空间是凸；甜甜圈（有洞）、月牙（有凹）、离散点集（不连通）不是凸。
\end{notebox}

\begin{center}
\begin{tikzpicture}[scale=0.80]
  % (a) 凸集正例：实心圆盘
  \begin{scope}
    \fill[gray!30] (0,0) circle (1.2);
    \draw[thick] (0,0) circle (1.2);
    \node[font=\footnotesize] at (0,-1.45) {凸集：实心圆盘};
  \end{scope}
  % (b) 凸集反例：圆环
  \begin{scope}[shift={(5.2,0)}]
    \fill[gray!30, even odd rule] (0,0) circle (1.2) (0,0) circle (0.45);
    \draw[thick] (0,0) circle (1.2);
    \draw[thick] (0,0) circle (0.45);
    \node[font=\footnotesize] at (0,-1.45) {非凸：圆环（有洞）};
  \end{scope}
  % (c) 非凸：月牙形
  \begin{scope}[shift={(0,-3.9)}]
    \fill[gray!30] (-0.3,0) circle (1.2);
    \fill[white] (0.45,0) circle (1.05);
    \draw[thick] (-0.3,0) circle (1.2);
    \draw[thick, dashed] (0.45,0) circle (1.05);
    \node[font=\footnotesize] at (0,-1.45) {非凸：月牙形（有凹坑）};
  \end{scope}
  % (d) 非凸：十字形
  \begin{scope}[shift={(5.2,-3.9)}]
    \fill[gray!30] (-1.2,-0.20) rectangle (1.2,0.20);
    \fill[gray!30] (-0.20,-1.2) rectangle (0.20,1.2);
    \draw[thick] (-1.2,-0.20) -- (1.2,-0.20) -- (1.2,0.20) -- (-1.2,0.20) -- cycle;
    \draw[thick] (-0.20,-1.2) -- (-0.20,1.2) -- (0.20,1.2) -- (0.20,-1.2) -- cycle;
    \node[font=\footnotesize] at (0,-1.45) {非凸：十字形（不实心）};
  \end{scope}
\end{tikzpicture}
\vspace{2pt}
{\small 图 5.1：凸集与非凸集示例\par}
\end{center}

\begin{defbox}
\textbf{定义 5.2（凸组合与凸包）}

点 $x_1,\dots,x_k$ 的\emph{凸组合}：$\sum_{i=1}^k \lambda_i x_i$，其中 $\lambda_i \ge 0$，$\sum \lambda_i = 1$。

集合 $S$ 的\emph{凸包} $\text{conv}(S)$：包含 $S$ 的最小凸集，等于 $S$ 中所有点的所有凸组合的集合。
\end{defbox}

\medskip
PPT第3.3页用图形展示凸包的构造过程：两点凸组合 $=$ 线段；三点凸组合 $=$ 三角形区域（含内部）。凸包就是把离散点集“填满”成包含它们的最小凸集。
\medskip
\refslide{3凸分析}{5}{3.3}
\medskip

\textbf{例3}：$S = \{(0,0),(1,0),(0,1)\}$，则 $\text{conv}(S)$ 是以此三点为顶点的直角三角形（含内部），即 $\{(x_1,x_2) \mid x_1+x_2 \le 1,\ x_1 \ge 0,\ x_2 \ge 0\}$。

\begin{defbox}
\textbf{定义 5.3（多面体）}

\emph{多面体}（polyhedral set）是有限个闭半空间的交：$\{x \mid Ax \le b\}$。LP的可行域 $\{x \mid Ax = b,\ x \ge 0\}$ 也是多面体——$Ax=b$ 等价于 $Ax\le b$ 且 $-Ax\le -b$，$x\ge0$ 等价于 $-x_i\le 0$。
\end{defbox}

\medskip
PPT第3.4页验证了凸集的常见例子（半空间、超平面、多面体、球等）；PPT第3.5页展示多面体的几何定义——有限个半空间的交集形成凸多面体。
\medskip
\refslide{3凸分析}{8}{3.4}
\medskip
\refslide{3凸分析}{9}{3.5}
\medskip

\begin{defbox}
\textbf{定义 5.4（极点）}

凸集 $C$ 中的点 $x$ 称为\emph{极点}（extreme point），如果它不能表示为 $C$ 中两个不同点的严格凸组合。即：
{\small
\[
x = \lambda y + (1-\lambda)z,\ \lambda\in(0,1),\ y,z\in C \implies x=y=z
\]
}
极点是凸集的“角”——LP的基可行解（BFS）就是可行域多面体的极点。
\end{defbox}

\medskip
PPT第3.6页给出极点的严格定义：$x$ 是极点当且仅当 $x=\lambda x_1+(1-\lambda)x_2$，$\lambda\in(0,1)$，$x_1,x_2\in S$ 推出 $x=x_1=x_2$。换言之，极点不能写成两个不同点的凸组合。PPT还指出：有界闭凸集中任一点都可表成极点的凸组合。
\medskip
\refslide{3凸分析}{11}{3.6}
\medskip

\textbf{例4（极点的几何直觉）}：
\begin{itemize}
  \item 三角形区域的极点 = 三个顶点
  \item 正方形/长方体的极点 = 四个/八个顶点
  \item 圆盘 $\{x \mid \|x\|_2 \le 1\}$ 的极点 = \emph{整个圆周}（无穷多个极点！）
  \item LP可行域 $\{x \mid Ax=b,\ x\ge0\}$ 的极点 = 基可行解（BFS）
\end{itemize}

\begin{notebox}
\textbf{极点与单纯形法的关系}：LP的最优解一定能在极点（BFS）处取得。单纯形法就是在极点之间移动，每次迭代从一个BFS跳到目标函数值更优的相邻BFS。这是第2章单纯形法的理论基础。
\end{notebox}

\begin{center}
\begin{tikzpicture}[scale=1.0]
  % 定义凸多边形顶点（逆时针）
  \coordinate (A) at (0,0);
  \coordinate (B) at (3.2,0.6);
  \coordinate (C) at (3.6,2.3);
  \coordinate (D) at (2.2,3.6);
  \coordinate (E) at (0.3,2.6);
  % 绘制填充多边形
  \draw[thick, fill=gray!15] (A) -- (B) -- (C) -- (D) -- (E) -- cycle;
  % 标记极点（顶点）
  \fill (A) circle (2.2pt) node[below] {$x^1$};
  \fill (B) circle (2.2pt) node[below right] {$x^2$};
  \fill (C) circle (2.2pt) node[right] {$x^3$};
  \fill (D) circle (2.2pt) node[above] {$x^4$};
  \fill (E) circle (2.2pt) node[above left] {$x^5$};
  % 内部点 x（两个极点的凸组合）
  \coordinate (P) at (1.9,2.0);
  \fill (P) circle (2.0pt) node[below] {$x$};
  % 虚线连接到两个极点
  \draw[dashed, gray, thin] (P) -- (B);
  \draw[dashed, gray, thin] (P) -- (D);
  % 标注
  \node[font=\footnotesize] at (2.4,1.5) {$x=\lambda x^2+(1-\lambda)x^4$};
  \node[font=\footnotesize] at (1.9,3.9) {$\bullet$ 极点（顶点）};
\end{tikzpicture}
\vspace{2pt}
{\small 图 5.3：极点示意图 —— 多边形顶点为极点，内部点 $x$ 可表为两个极点的凸组合\par}
\end{center}

\begin{defbox}
\textbf{定理 5.1（表示定理，选读）}

设多面体 $P=\{x \mid Ax \le b\} \neq \emptyset$ 且 $\text{rank}(A)=n$。记 $P$ 的极点集为 $\{x^k\}_{k\in K}$（非空有限），极方向集为 $\{d^j\}_{j\in J}$（$J$ 可为空）。则：
\[
x \in P \;\Longleftrightarrow\; x = \sum_{k\in K}\lambda_k x^k + \sum_{j\in J}\mu_j d^j
\]
其中 $\lambda_k \ge 0,\ \sum\lambda_k = 1,\ \mu_j \ge 0$。此外，$J=\emptyset \Longleftrightarrow P$ 有界（称为多胞形）。
\end{defbox}

\medskip
PPT第3.8页给出表示定理的完整陈述和公式——多面体中任意点都可以分解为极点凸组合与极方向非负组合的和。PPT第3.9页给出重要 \emph{推论}：标准形多面体 $\{x \mid Ax=b,\ x\ge 0\}$ 非空则必有极点。这保证了单纯形法的理论基础——只需在极点（BFS）之间迭代搜索。
\medskip
\refslide{3凸分析}{14}{3.8}
\medskip
\refslide{3凸分析}{15}{3.9}
\medskip

\begin{notebox}
\textbf{考试定位}：表示定理属选读内容，25年未直接出题。核心记住\emph{推论}——标准形LP可行域非空则BFS存在，单纯形法完备。
\end{notebox}

% ============================================================
\subsection{5.3 凸函数}

\begin{defbox}
\textbf{定义 5.5（凸函数）}

设 $C \subseteq \mathbb{R}^n$ 为非空凸集，$f: C \to \mathbb{R}$。$f$ 称为\emph{凸函数}，若 $\forall x,y\in C,\ \lambda\in(0,1)$：
\[
f(\lambda x + (1-\lambda) y) \le \lambda f(x) + (1-\lambda) f(y)
\]
若 $x\neq y$ 时严格不等式成立，称为\emph{严格凸}。若 $-f$ 是凸函数，则 $f$ 为\emph{凹函数}；若 $-f$ 是严格凸函数，则 $f$ 为\emph{严格凹函数}。
\end{defbox}

\medskip
PPT第3.10页给出凸函数形式定义与“弦在曲线上方”的几何图示。
\medskip
\refslide{3凸分析}{17}{3.10}
\medskip

\begin{notebox}
\textbf{几何解释（弦在曲线上方）}：任取定义域内两点 $x,y$，连接 $(x,f(x))$ 与 $(y,f(y))$ 的直线段 \emph{整个在函数图像上方}。凹函数则反过来——弦在图像下方。仿射函数（$a^Tx+b$）两边都等号，既凸又凹。
\end{notebox}

\begin{center}
\begin{tikzpicture}[x=1.2cm, y=0.55cm, >=stealth]
  % 坐标轴
  \draw[->,thick] (-0.15,0) -- (3.7,0) node[right] {$x$};
  \draw[->,thick] (0,-0.15) -- (0,10.8) node[above] {$y$};
  % 曲线 y = x^2
  \draw[thick, domain=0:3.2, samples=80] plot (\x, {\x*\x});
  \node[font=\small] at (3.5,10.0) {$y=x^2$};
  % 点 A(1,1) 和 B(3,9)
  \fill (1,1) circle (2.0pt) node[below left] {$A(1,1)$};
  \fill (3,9) circle (2.0pt) node[above right] {$B(3,9)$};
  % 弦（虚线）
  \draw[dashed, thick] (1,1) -- (3,9);
  % 标注“弦在曲线上方”
  \draw[->, gray] (2.05,5.6) -- (2.5,7.0) node[right, font=\small] {弦在曲线上方};
\end{tikzpicture}
\vspace{2pt}
{\small 图 5.2：凸函数的几何解释（弦在曲线上方）\par}
\end{center}

\textbf{例5（常见凸函数速查）}：

{\small
\begin{tabular}{@{}lll@{}}
\toprule
\textbf{函数} & \textbf{定义域} & \textbf{凸性} \\
\midrule
$a^T x+b$（仿射） & $\mathbb{R}^n$ & 既凸又凹 \\
$x^2$    & $\mathbb{R}$   & 严格凸 \\
$e^{ax}$ & $\mathbb{R}$   & 凸 \\
$|x|$    & $\mathbb{R}$   & 凸（$x=0$不可微，但仍凸！） \\
$x\ln x$ & $(0,\infty)$   & 凸（负熵） \\
$\|x\|$（任意范数）& $\mathbb{R}^n$ & 凸（三角不等式+齐次性）\\
$x_1^2+x_2^2$ & $\mathbb{R}^2$ & 严格凸 \\
$x_1^2-x_2^2$ & $\mathbb{R}^2$ & \textbf{非凸非凹}（鞍面） \\
$x^3$    & $\mathbb{R}$   & \textbf{非凸非凹}（$x>0$凸，$x<0$凹）\\
\bottomrule
\end{tabular}
}

\bigskip

% ============================================================
\subsection{5.4 凸函数的判别方法}

判断一个函数是否凸，有\emph{三层递进方法}。三者的\emph{因果关系}是逐步增加可微性假设，以此换取更强的计算可操作性：

\begin{center}
\fbox{%
\begin{minipage}{0.90\textwidth}
\small
\textbf{方法I：定义法}（无需可微假设，适合理论证明与简单函数）\\[2pt]
$\Big\Downarrow$ \quad 加上 $f$ 一阶可微\\[2pt]
\textbf{方法II：一阶条件/梯度不等式}（适合理论推导，但仍需“$\forall x,y$”）\\[2pt]
$\Big\Downarrow$ \quad 加上 $f$ 二阶可微\\[2pt]
\textbf{方法III：二阶条件/Hessian半正定}（\textbf{最适合具体计算}——顺序主子式法系统可操作）
\end{minipage}%
}
\end{center}

\textbf{为什么二阶判别最实用？} 定义法须证不等式对\emph{任意两点}成立；一阶条件虽简化了形式，仍需在“所有 $x,y\in S$”上验证。二阶条件将凸性判定转化为\emph{矩阵半正定性的代数判定}——计算Hessian矩阵，检查顺序主子式。对给定的具体函数，只需微积分 $+$ 初等代数，无需处理全称量词。这是考试中最常用的判别手段。

\begin{defbox}
\textbf{方法 5.1（判别法一：定义法）}

直接验证 $f(\lambda x+(1-\lambda)y) \le \lambda f(x)+(1-\lambda)f(y)$。无需可微性，仅适用于简单函数或理论证明。
\end{defbox}

\textbf{例6（定义法验证 $x^2$ 凸）}：$f(x)=x^2$。
{\small
\begin{align*}
& f(\lambda x+(1-\lambda)y) - [\lambda f(x)+(1-\lambda)f(y)] \\
&\quad = (\lambda x+(1-\lambda)y)^2 - \lambda x^2 - (1-\lambda)y^2 \\
&\quad = \lambda^2 x^2 + 2\lambda(1-\lambda)xy \\
&\qquad + (1-\lambda)^2 y^2 - \lambda x^2 - (1-\lambda)y^2 \\
&\quad = -\lambda(1-\lambda)(x^2 - 2xy + y^2) \\
&\quad = -\lambda(1-\lambda)(x-y)^2 \le 0 \quad\checkmark
\end{align*}
}

\textbf{例7（定义法验证范数凸）}：$f(x)=\|x\|$（不可微，只能用定义法）。
{\small
\[
\|\lambda x+(1-\lambda)y\| \le \|\lambda x\| + \|(1-\lambda)y\|
= \lambda\|x\| + (1-\lambda)\|y\| \quad\checkmark
\]
}
关键是三角不等式和范数的齐次性。

\begin{defbox}
\textbf{方法 5.2（判别法二：一阶条件/梯度不等式）}

设 $f$ 在开凸集 $S$ 上可微。则：
{\small
\[
f \text{ 凸 } \;\Longleftrightarrow\;
f(y) \ge f(x) + \nabla f(x)^T(y-x),\quad \forall x,y\in S
\]
}
几何含义：函数图像在任意点处的切线（一阶Taylor展开）始终在函数图像\emph{下方}。

（严格凸对应的不等号为严格 $>$，$\forall x\neq y$。）
\end{defbox}

\medskip
PPT第3.12页给出定理7：可微函数凸的充要条件是一阶梯度不等式对所有点成立。
\medskip
\refslide{3凸分析}{30}{3.12}
\medskip

\textbf{例8（一阶条件验证 $e^x$ 凸）}：$\nabla f(x)=e^x$。需证 $e^y \ge e^x + e^x(y-x)$，即 $e^{y-x} \ge 1+(y-x)$。令 $t=y-x$，需证 $e^t \ge 1+t$（$\forall t\in\mathbb{R}$）。这是指数函数的基本不等式，由 $e^t$ 的凸性可知成立。$\checkmark$

\begin{defbox}
\textbf{方法 5.3（判别法三：二阶条件/Hessian半正定）\quad $\star$ 考试最常用}

设 $f$ 在开凸集 $S$ 上二次可微。则：
\[
f \text{ 凸 } \;\Longleftrightarrow\; \nabla^2 f(x) \succeq 0,\quad \forall x\in S
\]
（$\nabla^2 f(x) \succ 0$ 对所有 $x$ 成立 $\Longrightarrow$ $f$ 严格凸，但反过来不成立！）

对于二次函数 $f(x)=\frac12 x^T Q x + c^T x$：$\nabla^2 f(x) \equiv Q$（常数矩阵），$f$ 凸 $\Longleftrightarrow$ $Q \succeq 0$。
\end{defbox}

\medskip
PPT第3.13页给出定理9：二次可微函数凸的充要条件是Hessian矩阵处处半正定；严格凸的充分条件是Hessian处处正定。这是考试中判别凸性最常用的定理。
\medskip
\refslide{3凸分析}{34}{3.13}
\medskip

\begin{notebox}
\textbf{顺序主子式法——Hessian正/半正定的代数判据}

对 $n\times n$ 实对称矩阵 $H$：（考试常见 $n=2$）
\begin{itemize}
  \item $H \succ 0$（正定） $\Longleftrightarrow$ 所有顺序主子式 $> 0$
  \item $H \succeq 0$（半正定）的判断需额外小心：顺序主子式 $\ge 0$ 是必要条件但不充分。但对常见的 $2\times2$ 情形：$\Delta_1 \ge 0$，$\Delta_2 \ge 0$ 可结合矩阵结构判断。
  \item $2\times2$ 矩阵快速判据：$H=\begin{pmatrix}a&b\\b&c\end{pmatrix}$。$H\succ0 \Longleftrightarrow a>0,\ ac-b^2>0$。$H\succeq0 \Longleftrightarrow a\ge0,\ c\ge0,\ ac-b^2\ge0$。
\end{itemize}
\end{notebox}

\textbf{例9（Hessian法练习——四种典型情形）}：

\textbf{(a) 正定（严格凸）}：$f(x_1,x_2) = 2x_1^2 + 3x_2^2$。
\[
\nabla f = \begin{pmatrix}4x_1 \\ 6x_2\end{pmatrix},\qquad
\nabla^2 f = \begin{pmatrix}4 & 0 \\ 0 & 6\end{pmatrix}
\]
$\Delta_1 = 4 > 0$，$\Delta_2 = 24 > 0$。$\nabla^2 f \succ 0$ $\Rightarrow$ $f$ 严格凸。

\textbf{(b) 正定（严格凸，带交叉项）}：$f(x_1,x_2) = x_1^2 + 2x_1x_2 + 3x_2^2$。
\[
\nabla^2 f = \begin{pmatrix}2 & 2 \\ 2 & 6\end{pmatrix}
\]
$\Delta_1 = 2 > 0$，$\Delta_2 = 2\cdot6-2\cdot2 = 8 > 0$。$\nabla^2 f \succ 0$ $\Rightarrow$ 严格凸。

\textbf{(c) 半正定（凸但不严格凸）}：$f(x_1,x_2) = x_1^2 + 4x_1x_2 + 4x_2^2 = (x_1+2x_2)^2$。
\[
\nabla^2 f = \begin{pmatrix}2 & 4 \\ 4 & 8\end{pmatrix}
\]
$\Delta_1 = 2 > 0$，$\Delta_2 = 2\cdot8-4\cdot4 = 0$。$\nabla^2 f \succeq 0$ 但非正定 $\Rightarrow$ $f$ 凸但不严格凸。函数沿方向 $d=(-2,1)^T$ 恒为零——在该方向上“平坦”。

\textbf{(d) 不定（非凸非凹）}：$f(x_1,x_2) = x_1^2 - 2x_1x_2 + \frac12 x_2^2$。
\[
\nabla^2 f = \begin{pmatrix}2 & -2 \\ -2 & 1\end{pmatrix}
\]
$\Delta_1 = 2 > 0$，$\Delta_2 = 2\cdot1-(-2)^2 = -2 < 0$。$\nabla^2 f$ 不定 $\Rightarrow$ $f$ 非凸非凹（鞍面）。

\begin{notebox}
\textbf{严格凸 $\neq$ Hessian处处正定（重要陷阱！）}

$\nabla^2 f \succ 0$（处处正定）是严格凸的\emph{充分但不必要条件}。反例：$f(x)=x^4$（$\mathbb{R}$上），$f''(x)=12x^2 \ge 0$，$f''(0)=0$（非正定），但 $f$ 仍为严格凸函数。考试中若 Hessian 处处正定则直接下结论“严格凸”；若 Hessian 半正定但非正定，仍需结合其他方法确认是否严格凸。
\end{notebox}

% ============================================================
\subsection{5.5 凸规划}

\begin{defbox}
\textbf{定义 5.6（凸规划）}

最优化问题 $\min\{f(x) \mid x \in S\}$ 称为\emph{凸规划}，如果：
\begin{itemize}
  \item 可行域 $S$ 是凸集
  \item 目标函数 $f$ 是凸函数（对 $\min$ 问题；对 $\max$ 则要求 $f$ 凹）
\end{itemize}
两个条件\emph{缺一不可}。
\end{defbox}

\begin{defbox}
\textbf{定理 5.2（凸规划的核心性质）}

对于凸规划 $\min\{f(x) \mid x \in S\}$：
\begin{enumerate}
  \item 任意\emph{局部}极小点都是\emph{全局}极小点
  \item 若 $f$ 严格凸，则全局极小点\emph{唯一}
  \item 上述结论\emph{不需要}任何约束规格——仅需 $f$ 凸且 $S$ 凸
  \item 若约束规格成立，KKT条件（将在第6章详细介绍）升级为全局最优性的\emph{充要条件}
  \item 全局极小点的集合是凸集
\end{enumerate}
\end{defbox}

\medskip
PPT第3.14页给出定理6：凸函数在凸集上的局部极小点一定是全局极小点，且极小点的集合为凸集。这是凸规划第一个核心性质的形式化陈述。
\medskip
\refslide{3凸分析}{27}{3.14}
\medskip

\begin{notebox}
\textbf{这为什么重要？} 第6章的KKT条件只是\emph{一阶必要条件}——满足KKT的点不一定是全局最优。但一旦确认问题是凸规划，KKT点自动升级为全局最优解。这是连接第5章（凸性）和第6章（KKT）的关键桥梁。

\textbf{判定流程}：目标函数凸？$\checkmark$ 可行域凸？$\checkmark$ $\Rightarrow$ 凸规划 $\Rightarrow$ KKT点=全局最优解。
\end{notebox}

% ============================================================
\subsection{5.6 凸函数的保凸运算}

记住以下运算规则，可在不计算Hessian的情况下快速判断组合函数的凸性：

{\small
\begin{tabular}{@{}ll@{}}
\toprule
\textbf{运算} & \textbf{条件与结论} \\
\midrule
非负加权和  & $f,g$ 凸，$\alpha,\beta \ge 0$ $\Rightarrow$ $\alpha f+\beta g$ 凸 \\
与仿射复合  & $f$ 凸 $\Rightarrow$ $f(Ax+b)$ 凸 \\
逐点最大值  & $f_1,\dots,f_m$ 凸 $\Rightarrow$ $\max_i f_i(x)$ 凸 \\
逐点上确界  & $\forall y\in A$，$f(\cdot,y)$ 凸 $\Rightarrow$ $\sup_{y\in A} f(x,y)$ 凸 \\
标量复合(1) & $g$ 凸，$h$ 凸且单调不减 $\Rightarrow$ $h(g(x))$ 凸 \\
标量复合(2) & $g$ 凹，$h$ 凸且单调不增 $\Rightarrow$ $h(g(x))$ 凸 \\
\bottomrule
\end{tabular}
}

\bigskip

\begin{notebox}
\textbf{实战技巧}：看到 $\max$ 结构立刻想到保凸运算。例如 $f(x)=\max\{x_1^2+x_2^2,\; 2x_1+3x_2\}$——$x_1^2+x_2^2$ 凸（$\nabla^2 f \succ 0$），$2x_1+3x_2$ 凸（仿射），逐点最大值保凸 $\Rightarrow$ $f$ 凸。无需算Hessian！
\end{notebox}

% ============================================================
\subsection{5.7 自编教学例题}

\subsubsection{题1：凸集判定}

\textbf{判定以下集合是否为凸集，并简述理由：}

(1) $S_1 = \{(x_1,x_2) \mid x_1^2 + x_2^2 \le 4,\ x_1 \ge 0\}$（右半圆盘）

(2) $S_2 = \{(x_1,x_2) \mid |x_1| + |x_2| \le 1\}$（$\ell_1$ 范数球/菱形）

(3) $S_3 = \{(x_1,x_2) \mid x_1^2 - x_2 \le 0\}$（抛物线 $x_2=x_1^2$ 的上方区域）

\medskip
\textbf{解：}

(1) 右半圆盘 $=$ 圆盘 $\cap$ 右半平面。两个凸集的交是凸集（凸集的性质）。$\checkmark$ 凸。

(2) $\ell_1$ 范数球 $=\{x \mid \|x\|_1 \le 1\}$。范数是凸函数，凸函数的下水平集是凸集。几何上是以 $(\pm1,0),(0,\pm1)$ 为顶点的菱形。$\checkmark$ 凸。

(3) $S_3 = \{x \mid x_2 \ge x_1^2\}$。$g(x)=x_1^2-x_2$ 是凸函数（$\nabla^2 g=\text{diag}(2,0)\succeq0$）。凸函数的下水平集是凸集。$\checkmark$ 凸。几何上：抛物线将平面分成“内部”（凹）和“外部”（凸），$S_3$ 是外部区域。

\subsubsection{题2：凸函数判定（Hessian法 + 顺序主子式）}

\textbf{判定以下二元函数的凸性：}

(1) $f(x_1,x_2) = 3x_1^2 + 2x_2^2 + 2x_1x_2$

(2) $f(x_1,x_2) = x_1^2 + x_2^4$

(3) $f(x_1,x_2) = -x_1^2 - 2x_2^2$\quad（注意负号！）

\medskip
\textbf{解：}

(1) $\nabla^2 f = \begin{pmatrix}6&2\\2&4\end{pmatrix}$。$\Delta_1=6>0$，$\Delta_2=6\cdot4-2\cdot2=20>0$。$\nabla^2 f \succ 0$ $\Rightarrow$ \emph{严格凸}。

(2) $\nabla^2 f = \begin{pmatrix}2&0\\0&12x_2^2\end{pmatrix}$。$\Delta_1=2>0$，$\Delta_2=24x_2^2 \ge 0$（对任意 $x$）。$\nabla^2 f \succeq 0$ 恒成立 $\Rightarrow$ $f$ 凸。注意 $x_2=0$ 处 Hessian 仅为半正定（非正定），但 $f$ 实际上仍是严格凸——这再次说明Hessian正定只是严格凸的充分条件。

(3) $\nabla^2 f = \begin{pmatrix}-2&0\\0&-4\end{pmatrix}$。$\Delta_1=-2<0$，$\Delta_2=8>0$。$\nabla^2 f \prec 0$（负定）$\Rightarrow$ $-f$ 严格凸 $\Rightarrow$ $f$ \emph{严格凹}。

\subsubsection{题3：凸规划实例——局部=全局的威力}

\textbf{考虑问题}：
\[
\begin{aligned}
\min\ & f(x_1,x_2) = x_1^2 + x_2^2 \\
\text{s.t.}\ & x_1 + x_2 \ge 1,\ x_1 \ge 0,\ x_2 \ge 0
\end{aligned}
\]

(1) 验证这是凸规划。\quad
(2) 求出全局最优解。

\medskip
\textbf{解：}

(1) 可行域 $S = \{x \mid x_1+x_2 \ge 1,\ x_1 \ge 0,\ x_2 \ge 0\}$ 是三个半空间的交 $\Rightarrow$ 多面体，凸。目标函数 $\nabla^2 f = \begin{pmatrix}2&0\\0&2\end{pmatrix} \succ 0$ $\Rightarrow$ 严格凸。$\therefore$ 这是凸规划（且 $f$ 严格凸，解唯一）。

(2) 几何解释：原点 $(0,0)$ 不可行（$0+0\ngeq1$）。全局最优解是可行域中离原点最近的点——即原点在半空间 $x_1+x_2\ge1$（正象限内）的投影。由对称性 $x_1=x_2$，结合 $x_1+x_2=1$ 得 $x_1^*=x_2^*=\frac12$。$f(x^*)=\frac12$。

KKT验证（凸规划下KKT充要）：构造 Lagrange 函数 $L = x_1^2+x_2^2 + \lambda(1-x_1-x_2) - \mu_1 x_1 - \mu_2 x_2$。
\[
\begin{aligned}
\frac{\partial L}{\partial x_1} &= 2x_1 - \lambda - \mu_1 = 0 \\
\frac{\partial L}{\partial x_2} &= 2x_2 - \lambda - \mu_2 = 0
\end{aligned}
\]
互补松弛：$\lambda(1-x_1-x_2)=0$，$\mu_1 x_1=0$，$\mu_2 x_2=0$。最优处 $x_1>0,x_2>0$ $\Rightarrow$ $\mu_1=\mu_2=0$。约束 $x_1+x_2\ge1$ 在最优处积极 $\Rightarrow$ $x_1+x_2=1$。代入得 $2x_1=\lambda=2x_2$ $\Rightarrow$ $x_1=x_2=\frac12$，$\lambda=1$。因为是凸规划，此KKT点即全局最优。

\subsubsection{题4：从第6章真题提取凸性片段}

25年真题第六题：
\[
\min\; x_1^2 + x_2^2,\quad \text{s.t.}\ (x_1-1)^3 - x_2^2 = 0
\]

\textbf{先看目标函数}：$f(x)=x_1^2+x_2^2$，$\nabla^2 f = \begin{pmatrix}2&0\\0&2\end{pmatrix} \succ 0$ $\Rightarrow$ $f$ 是\emph{严格凸函数}。

\textbf{再看约束}：$g(x)=(x_1-1)^3-x_2^2=0$。考虑曲线上的两点：$x^A=(2,1)$（$g=1^3-1=0$）和 $x^B=(2,-1)$（$g=1^3-1=0$）。中点 $x^M=(2,0)$，$g(x^M)=1^3-0=1\neq0$，不在约束曲线上。约束集合\emph{非凸}。

\textbf{结论}：目标凸 $+$ 约束非凸 $\Rightarrow$ \emph{不是凸规划}。KKT条件仅为\emph{必要条件}——这正是该题有FJ点 $(1,0)$ 但无KKT点的根本原因。凸规划判定必须两个条件同时满足，缺一不可！

\begin{notebox}
\textbf{判定凸规划的完整流程（考试操作步骤）}

Step 1：检查目标函数——求Hessian、算顺序主子式，判断是否半正定/正定。\\
Step 2：检查可行域——是否可写为 $\{x\mid g_i(x)\le0,\ h_j(x)=0\}$ 形式，每个 $g_i$ 是否凸（不等号 $\le$）、每个 $h_j$ 是否仿射（等号）。\\
Step 3：两个条件都满足 $\Rightarrow$ 凸规划 $\Rightarrow$ 局部$=$全局，KKT充要。\\
\hspace{1.5em}任一不满足 $\Rightarrow$ 非凸规划 $\Rightarrow$ KKT仅为必要条件。
\end{notebox}

% ============================================================
\clearpage\subsection{本章速查}

\begin{itemize}
  \item \textbf{凸集}：线段全在集合内。多面体（$Ax\le b$）、半空间、超平面、范数球都凸
  \item \textbf{凸集反例检验}：有洞（圆环）、有凹（月牙）、离散（$\mathbb{Z}^n$）、不连通状
  \item \textbf{凸函数判别三层次}：定义法 $\to$ 一阶梯度不等式 $\to$ 二阶Hessian半正定
  \item \textbf{二阶判别最实用}：算Hessian $\to$ 顺序主子式 $\to$ 正定/半正定/不定/负定
  \item \textbf{$2\times2$ Hessian快判}：$a>0,\ ac-b^2>0$ 为正定；$ac-b^2<0$ 为不定
  \item \textbf{凸规划}：凸可行域 $+$ 凸目标（对 $\min$）$\Rightarrow$ 局部 $=$ 全局，KKT充要
  \item \textbf{凸规划判定}：目标凸性和可行域凸性\emph{缺一不可}，必须分别验证
  \item \textbf{严格凸} $\Rightarrow$ 最优解唯一（但 Hessian 正定仅仅是充分条件）
  \item \textbf{保凸运算}：非负加权和、与仿射复合、逐点最大值均保持凸性
  \item \textbf{常见凸函数}：仿射、$x^2$、$e^x$、$|x|$、范数、$x\ln x$（负熵）
  \item \textbf{常见非凸/非凹}：$x_1^2-x_2^2$（鞍面）、$x^3$（$\mathbb{R}$上）、$\sin x$
\end{itemize}

% ============================================================
%  第6章：最优性条件（FJ条件与KKT条件）
% ============================================================
\section{第6章 \quad 最优性条件}

\subsection{6.1 从无约束到有约束}

无约束优化 $\min f(x)$ 的最优性条件很简单：$\nabla f(x^*) = 0$（一阶必要条件），外加 Hessian 半正定（二阶必要条件）。

加上约束后，梯度为零不再适用——最优解可能在约束边界上，梯度不需要为零。我们需要一个新的工具来描述"在约束允许的范围内，目标函数不能进一步改进"。

\subsection{6.2 等式约束与 Lagrange 乘子法}

先看最简单的情形：仅有等式约束。

\begin{defbox}
\textbf{定理 6.1（Lagrange 乘子法）}

考虑 $\min f(x)$，s.t.\ $h_i(x)=0$，$i=1,\dots,m$。定义 Lagrange 函数：
\[
L(x, \lambda) = f(x) + \sum_{i=1}^m \lambda_i h_i(x)
\]
若 $x^*$ 是局部极小点且 $\nabla h_i(x^*)$ 线性无关（LICQ 成立），则存在 $\lambda_i^*$ 使：
\[
\nabla_x L(x^*, \lambda^*) = \nabla f(x^*) + \sum_{i=1}^m \lambda_i^* \nabla h_i(x^*) = 0
\]
即 $-\nabla f(x^*) = \sum \lambda_i^* \nabla h_i(x^*)$：负梯度在约束法向量张成的空间中。
\end{defbox}

\medskip
\refslide{7最优性条件}{3}{7.1}
\medskip

\begin{notebox}
\textbf{几何解释}：在约束曲面上的局部极小点处，目标函数的负梯度一定可以表示为约束梯度的线性组合。如果梯度在切空间方向上有分量，沿着那个分量移动就能在约束曲面上进一步下降。
\end{notebox}

\subsection{6.3 不等式约束与 FJ 条件}

考虑一般约束问题：
\[
\min\ f(x),\quad \text{s.t.}\ g_i(x) \le 0\ (i=1,\dots,m),\ h_j(x)=0\ (j=1,\dots,p)
\]

\begin{defbox}
\textbf{定理 6.2（Fritz John 条件，FJ 条件）}

设 $x^*$ 是局部极小点。则存在不全为零的乘子 $\lambda_0, \lambda_1,\dots,\lambda_m, \mu_1,\dots,\mu_p$ 满足：
\[
\begin{aligned}
& \lambda_0 \nabla f(x^*) + \sum_{i=1}^m \lambda_i \nabla g_i(x^*) + \sum_{j=1}^p \mu_j \nabla h_j(x^*) = 0 \\
& \lambda_i \ge 0 \ (i=0,1,\dots,m) \\
& \lambda_i g_i(x^*) = 0 \ (i=1,\dots,m) \quad \text{（互补松弛）}
\end{aligned}
\]
\end{defbox}

\medskip
\refslide{7最优性条件}{17}{7.9}
\medskip

\begin{notebox}
\textbf{FJ 条件的致命弱点}：$\lambda_0$ 可以为零。如果 $\lambda_0 = 0$，那么方程组中目标函数的梯度\emph{完全消失}——只剩约束梯度的线性组合为零。此时 FJ 条件给出的点可能与目标函数没有任何关系。这就是所谓的"退化"情况。
\end{notebox}

\subsection{6.4 KKT 条件——FJ 的强化版}

KKT 条件在 FJ 条件的基础上\emph{加上约束规格}（constraint qualification），保证 $\lambda_0 \neq 0$（常归一化为 1），从而目标函数的梯度信息不丢失。

\begin{defbox}
\textbf{定理 6.3（KKT 条件）}

设 $x^*$ 是局部极小点，且在 $x^*$ 处某约束规格成立（如 LICQ 或 Slater 条件）。则存在乘子 $\lambda_i \ge 0$（对 $g_i$）和 $\mu_j$（对 $h_j$），使：
\[
\begin{aligned}
& \nabla f(x^*) + \sum_{i=1}^m \lambda_i \nabla g_i(x^*) + \sum_{j=1}^p \mu_j \nabla h_j(x^*) = 0 \\
& \lambda_i g_i(x^*) = 0 \quad (i=1,\dots,m) \quad \text{（互补松弛）} \\
& g_i(x^*) \le 0,\ h_j(x^*) = 0 \quad \text{（原始可行性）} \\
& \lambda_i \ge 0 \quad \text{（对偶可行性）}
\end{aligned}
\]
(与 FJ 的核心区别：$\lambda_0$ 已被归一化为 1，且乘子 $\mu_j$ 不再有符号约束。)
\end{defbox}

\medskip
\refslide{7最优性条件}{19}{7.11}
\medskip
\refslide{7最优性条件}{20}{7.12}
\medskip

\begin{defbox}
\textbf{常见约束规格}

\begin{itemize}
  \item \textbf{LICQ（线性独立约束规格）}：积极约束的梯度向量集合 $\{\nabla g_i(x^*) : g_i(x^*)=0\} \cup \{\nabla h_j(x^*)\}$ 线性无关。最常用。
  \item \textbf{Slater 条件}（对凸问题）：存在严格可行点使 $g_i(x) < 0$。
  \item \textbf{线性约束自动满足}：若所有约束都是线性的，则约束规格自动成立，无需额外验证。
\end{itemize}
\end{defbox}

\medskip
\refslide{7最优性条件}{23}{7.15}
\medskip

\subsection{6.5 FJ vs KKT 对比}

\begin{defbox}
\textbf{表 6.1：FJ 与 KKT 的区别}

\[
\begin{array}{l|l|l}
& \text{FJ 条件} & \text{KKT 条件} \\ \hline
\lambda_0 & \ge 0\text{，可以为 }0 & \text{归一化为 }1\ (\lambda_0 \neq 0) \\
\text{约束规格} & \text{不需要} & \text{必须满足} \\
\text{适用性} & \text{所有局部极小点} & \text{仅满足约束规格的局部极小点} \\
\text{信息量} & \text{可能退化（$\lambda_0=0$）} & \text{始终包含梯度信息} \\
\text{考试地位} & \text{更一般的理论框架} & \text{实际求解的主要工具}
\end{array}
\]
\end{defbox}

\medskip
\refslide{7最优性条件}{26}{7.18}
\medskip

\subsection{6.6 真题示例：第六题(1)(2)——FJ 条件与 KKT}

\begin{notebox}
\textbf{【2025 真题·第六题】}
\[
\min\ x_1^2 + x_2^2,\quad \text{s.t.}\ (x_1-1)^3 - x_2^2 = 0
\]
(1) 找出所有满足 FJ 条件的可行点；(2) 证明不存在可行点满足 KKT 条件。
\end{notebox}

\medskip

\subsubsection{(1) FJ 条件求解}

目标函数梯度：$\nabla f(x) = (2x_1, 2x_2)^T$。约束梯度：$\nabla g(x) = (3(x_1-1)^2, -2x_2)^T$。

FJ 方程（存在不全为零的 $\lambda_0 \ge 0, \mu$）：
\[
\lambda_0 \begin{pmatrix} 2x_1 \\ 2x_2 \end{pmatrix}
+ \mu \begin{pmatrix} 3(x_1-1)^2 \\ -2x_2 \end{pmatrix}
= \begin{pmatrix} 0 \\ 0 \end{pmatrix}
\]
且 $g(x) = (x_1-1)^3 - x_2^2 = 0$。

\textbf{情形 1：$\lambda_0 = 0$}。则 $\mu \neq 0$（不能全是零）。方程化为：
\[
\mu \cdot 3(x_1-1)^2 = 0,\quad \mu \cdot (-2x_2) = 0
\]
$\mu \neq 0 \Rightarrow x_1 = 1,\ x_2 = 0$。代入约束：$(0)^3 - 0^2 = 0$，满足。得 FJ 点：$(1,0)$。

\textbf{情形 2：$\lambda_0 \neq 0$}。归一并消去 $\lambda_0$（设 $\lambda_0=1$，$\mu$ 重新尺度化）：
\[
\begin{pmatrix} 2x_1 \\ 2x_2 \end{pmatrix}
+ \mu \begin{pmatrix} 3(x_1-1)^2 \\ -2x_2 \end{pmatrix}
= 0
\]
从第二个方程：$2x_2 - 2\mu x_2 = 0 \Rightarrow 2(1-\mu)x_2 = 0 \Rightarrow \mu = 1$ 或 $x_2 = 0$。

若 $\mu=1$：$2x_1 + 3(x_1-1)^2 = 0$，此方程无实根。\\
若 $x_2=0$：从约束得 $x_1=1$。代入第一方程：$2\cdot 1 + \mu\cdot 3(0)^2 = 2 \neq 0$，矛盾。

故仅有唯一的 FJ 可行点：$(1,0)^T$。

\medskip
\refslide{7最优性条件}{17}{7.9}
\medskip

\subsubsection{(2) 证明 KKT 不存在}

KKT 要求 $\lambda_0 \neq 0$。由上分析，情形 2 中唯一可能的候选 $(1,0)$ 代入第一方程得 $2 \neq 0$，不是 KKT 点。故不存在 KKT 点。

\begin{notebox}
\textbf{这题揭示了什么？} FJ 条件找到了一个点 $(1,0)$，但它不是 KKT 点。原因是该点处约束梯度 $\nabla g(1,0) = (0,0)^T = 0$——约束梯度退化，不满足 LICQ。这是 FJ 条件"$\lambda_0=0$ 退化"的经典考试情景。
\end{notebox}

\medskip
\refslide{7最优性条件}{27}{7.19}
\medskip

\subsection{本章速查}

\begin{itemize}
  \item FJ：$\lambda_0 \nabla f + \sum\lambda_i\nabla g_i + \sum\mu_j\nabla h_j = 0$，$\lambda_0\ge 0$（可为 0），$\lambda_i\ge 0$，$\lambda_i g_i=0$
  \item KKT：同 FJ 但 $\lambda_0=1$（已归一化），需要约束规格
  \item 约束规格：LICQ（积极约束梯度线性无关）、Slater、线性约束自动满足
  \item 核心区别：FJ 允许 $\lambda_0=0$（目标信息消失），KKT 不允许
  \item 考试关键技巧：FJ 题分 $\lambda_0=0$ 和 $\lambda_0\neq 0$ 两种情形讨论
\end{itemize}

% ============================================================
%  第7章：非线性对偶
% ============================================================
\section{第7章 \quad 非线性对偶}

\subsection{7.1 从线性对偶到非线性对偶}

第3章讲的线性规划对偶，结论非常漂亮：强对偶定理成立，对偶间隙为零。到了非线性规划，事情变得复杂——强对偶不总是成立。

非线性对偶的核心工具是 \textbf{Lagrange 对偶函数}。思路：把约束"价格化"（乘子），问在没有约束的情况下，目标+惩罚的最坏情况是什么。

\subsection{7.2 Lagrange 对偶函数}

考虑一般非线性规划（原问题）：
\[
\min\ f(x),\quad \text{s.t.}\ g_i(x) \le 0\ (i=1,\dots,m),\ h_j(x)=0\ (j=1,\dots,p),\ x \in X
\]

\begin{defbox}
\textbf{定义 7.1（Lagrange 对偶函数）}

给定乘子 $\lambda \ge 0$（对不等式约束）和 $\mu$（对等式约束），定义 Lagrange 对偶函数：
\[
\theta(\lambda, \mu) = \inf_{x \in X} \left\{ f(x) + \sum_{i=1}^m \lambda_i g_i(x) + \sum_{j=1}^p \mu_j h_j(x) \right\}
\]
即在乘子给定的情况下，对 $x$ 取下确界。对偶函数提供了原问题最优值的下界。
\end{defbox}

\medskip
\refslide{7最优性条件}{34}{7.22}
\medskip
\refslide{7最优性条件}{35}{7.23}
\medskip

\begin{notebox}
\textbf{为什么要取下确界？} 原问题是带约束的 $\min$。取下确界相当于"放松了所有约束，但用乘子给予奖惩"——得到的值是原问题的一个下界。乘子不同，下界也不同。
\end{notebox}

\subsection{7.3 弱对偶定理（恒成立）}

\begin{defbox}
\textbf{定理 7.1（弱对偶定理）}

设 $x$ 是原问题的任意可行解，$(\lambda, \mu)$ 是任意对偶乘子（$\lambda \ge 0$）。则恒有：
\[
f(x) \ge \theta(\lambda, \mu)
\]
因此：
\[
\min\ f(x) \ge \max_{\lambda \ge 0,\ \mu}\ \theta(\lambda, \mu)
\]
即：原最优值 $\ge$ 对偶最优值。这个差值称为\emph{对偶间隙}。
\end{defbox}

\medskip
\refslide{7最优性条件}{36}{7.24}
\medskip

\textbf{弱对偶不需要任何条件}——不需要凸性，不需要约束规格，永远成立。

\subsection{7.4 强对偶定理（需要条件）}

\begin{defbox}
\textbf{定理 7.2（强对偶定理）}

若原问题为凸规划（$f$ 凸，$g_i$ 凸，$h_j$ 仿射）且 Slater 条件成立（存在严格可行点使 $g_i(x) < 0$），则：
\[
\min\ f(x) = \max_{\lambda \ge 0,\ \mu}\ \theta(\lambda, \mu)
\]
且对偶问题的最优解（乘子）就是原问题 KKT 条件中的乘子。
\end{defbox}

\medskip
\refslide{7最优性条件}{48}{7.32}
\medskip

\subsection{7.5 真题示例：第六题(3)(4)——非线性对偶与变量消元}

回到第六题：
\[
\min\ x_1^2 + x_2^2,\quad \text{s.t.}\ (x_1-1)^3 - x_2^2 = 0
\]

\subsubsection{(3) 变量消元法的正确性}

有人尝试：从约束解出 $x_2^2 = (x_1-1)^3$，代入目标函数得：
\[
\min\ x_1^2 + (x_1-1)^3
\]
这种做法正确吗？

\textbf{错误！} 遗漏了隐含条件。$x_2^2 = (x_1-1)^3$ 意味着左边 $x_2^2 \ge 0$，所以右边也必须 $\ge 0$：
\[
(x_1-1)^3 \ge 0 \quad\Rightarrow\quad x_1 \ge 1
\]
代入后的一元问题必须加上 $x_1 \ge 1$：
\[
\min\ x_1^2 + (x_1-1)^3,\quad x_1 \ge 1
\]
$g(x_1) = x_1^2 + (x_1-1)^3$，$g'(x_1) = 2x_1 + 3(x_1-1)^2 > 0$（当 $x_1 \ge 1$），故 $x_1=1$ 时取最小：$f^* = 1$，$x^* = (1, 0)^T$。若无 $x_1 \ge 1$，原表达式 $x_1^2 + (x_1-1)^3$ 当 $x_1 \to -\infty$ 时 $\to -\infty$，无下界——这就是遗漏隐含约束的后果。

\medskip
\refslide{7最优性条件}{51}{7.35}
\medskip

\subsubsection{(4) 对偶规划}

构造 Lagrange 对偶函数（仅一个等式约束，乘子为 $\mu$）：
\[
L(x, \mu) = x_1^2 + x_2^2 + \mu[(x_1-1)^3 - x_2^2]
\]

对偶函数 $\theta(\mu) = \inf_x L(x, \mu)$：
\begin{itemize}
  \item 若 $\mu = 0$：$L(x,0) = x_1^2 + x_2^2$，下确界为 $0$
  \item 若 $\mu \neq 0$：当 $x_2 \to \infty$ 时（暂忽略约束），$x_2^2 - \mu x_2^2 = (1-\mu)x_2^2$。若 $1-\mu < 0$（即 $\mu > 1$），$x_2^2$ 前面系数为负，$\to -\infty$；更仔细的分析表明对所有 $\mu \neq 0$，下确界都是 $-\infty$
\end{itemize}
故：
\[
\theta(\mu) = \begin{cases}
0, & \mu = 0 \\
-\infty, & \mu \neq 0
\end{cases}
\]

对偶规划：$\max\ \theta(\mu) = 0$。

\begin{notebox}
\textbf{注意}：原最优值 $f^* = 1$，对偶最优值 $= 0$。对偶间隙 $= 1$。为什么强对偶不成立？因为这个问题\emph{不是凸的}——约束 $(x_1-1)^3 - x_2^2 = 0$ 并非凸约束。弱对偶成立（$1 > 0$），强对偶不成立（间隙不为零）。
\end{notebox}

\medskip
\refslide{7最优性条件}{52}{7.36}
\medskip
\refslide{7最优性条件}{55}{7.39}
\medskip

\subsection{本章速查}

\begin{itemize}
  \item Lagrange 对偶函数：$\theta(\lambda,\mu) = \inf_{x\in X} \{f(x) + \sum\lambda_i g_i(x) + \sum\mu_j h_j(x)\}$
  \item 弱对偶定理：$f(x) \ge \theta(\lambda,\mu)$，永恒成立
  \item 强对偶定理：需凸规划 + 约束规格，原最优值 = 对偶最优值
  \item 对偶间隙：原-对偶最优值之差，非凸问题可能 $>0$
  \item 变量消元陷阱：必须保留隐含不等式约束（如 $x^2 \ge 0$ 带来的 $x_1 \ge 1$）
\end{itemize}

% ============================================================
%  第8章：无约束优化方法
% ============================================================
\section{第8章 \quad 无约束优化方法}

\subsection{8.1 下降算法的通用框架}

无约束优化 $\min f(x)$ 不能用单纯形表——我们没有约束提供的结构。但我们可以用迭代法：从某点出发，每一步选择一个\emph{下降方向}，沿该方向走\emph{合适的步长}。

\begin{defbox}
\textbf{下降算法通用格式}

给定初始点 $x^{(0)}$。对 $k = 0,1,2,\dots$：
\[
x^{(k+1)} = x^{(k)} + \lambda_k d^{(k)}
\]
其中：
\begin{itemize}
  \item $d^{(k)}$ 是\emph{搜索方向}，满足 $\nabla f(x^{(k)})^T d^{(k)} < 0$（下降方向）
  \item $\lambda_k > 0$ 是\emph{步长}，使得 $f(x^{(k+1)}) < f(x^{(k)})$
\end{itemize}
\end{defbox}

\medskip
\refslide{10使用导数的优化算法}{40}{10.1}
\medskip

不同的无约束优化方法，区别在于\emph{如何选方向}和\emph{如何定步长}。

\subsection{8.2 最速下降法（梯度法）}

\begin{defbox}
\textbf{最速下降法}

取 $d^{(k)} = -\nabla f(x^{(k)})$（负梯度方向）。这是局部下降最快的方向。步长通过精确线搜索确定：
\[
\lambda_k = \argmin_{\lambda \ge 0}\ f(x^{(k)} + \lambda d^{(k)})
\]
\end{defbox}

最速下降法实现简单，但收敛慢（呈现锯齿形）。在 Hessian 条件数大时尤为严重。它为理解更高级方法提供了基准。

\subsection{8.3 精确线搜索}

在选定方向 $d^{(k)}$ 后，需要确定步长 $\lambda_k$。精确线搜索求解一元极小化问题 $\min_{\lambda \ge 0}\ \phi(\lambda) = f(x^{(k)} + \lambda d^{(k)})$。

\textbf{做法}：将 $x^{(k)} + \lambda d^{(k)}$ 代入 $f$ 得 $\phi(\lambda)$，求导 $\phi'(\lambda) = 0$ 解出 $\lambda$。

\medskip
\refslide{9一维搜索}{51}{9.8}
\medskip
\refslide{9一维搜索}{52}{9.9}
\medskip

\subsection{8.4 共轭梯度法}

最速下降法的锯齿现象源于\emph{连续两步的搜索方向是正交的}（一碰就转 90°）。共轭梯度法通过巧妙的递推\emph{构造共轭方向}，使得对正定二次函数 $f(x) = \frac{1}{2}x^T Q x + c^T x$，$n$ 步内精确收敛。

\begin{defbox}
\textbf{FR（Fletcher-Reeves）共轭梯度法}

初始化：$d^{(0)} = -\nabla f(x^{(0)})$。

对于 $k \ge 1$：
\[
\beta_k = \frac{\| \nabla f^{(k)} \|^2}{\| \nabla f^{(k-1)} \|^2}
\]
\[
d^{(k)} = -\nabla f^{(k)} + \beta_k d^{(k-1)}
\]
方向 $d^{(k)}$ 与 $d^{(k-1)}$ 关于 Hessian $Q$ 共轭：$(d^{(k)})^T Q d^{(k-1)} = 0$。
\end{defbox}

\medskip
\refslide{10使用导数的优化算法}{46}{10.7}
\medskip

\begin{notebox}
\textbf{共轭梯度法为什么好？} 它不需要存储 Hessian 矩阵（存储量 $O(n)$ vs 牛顿法 $O(n^2)$），适合大规模问题；收敛速度远快于最速下降法；利用了过去方向的信息避免了锯齿。
\end{notebox}

\subsection{8.5 拟牛顿法}

\subsubsection{牛顿法的瓶颈}

牛顿法方向 $d^{(k)} = -[\nabla^2 f(x^{(k)})]^{-1} \nabla f(x^{(k)})$ 收敛极快（二次收敛），但需要：一、计算 Hessian 矩阵；二、求 Hessian 逆。两者计算量都很大。

\subsubsection{拟牛顿法的核心思想}

\textbf{不直接算 Hessian 逆，而是用迭代过程中观察到的梯度差来逐步"逼近"它}。

设 $s_{k-1} = x^{(k)} - x^{(k-1)}$（位移差），$y_{k-1} = \nabla f^{(k)} - \nabla f^{(k-1)}$（梯度差）。拟牛顿条件要求近似 Hessian 逆 $H_k$ 满足：
\[
H_k y_{k-1} = s_{k-1}
\]
这有明确的几何含义：$H_k$ 沿着方向 $y_{k-1}$ 的"曲率"应该匹配实际观察到的"曲率"。

\subsubsection{DFP 校正公式}

\begin{defbox}
\textbf{DFP（Davidon-Fletcher-Powell）公式}

从 $H_{k-1}$ 更新到 $H_k$ 的公式为：
\[
H_k = H_{k-1} + \underbrace{\frac{s_{k-1} s_{k-1}^T}{s_{k-1}^T y_{k-1}}}_{\text{正定修正}}
- \underbrace{\frac{H_{k-1} y_{k-1} y_{k-1}^T H_{k-1}}{y_{k-1}^T H_{k-1} y_{k-1}}}_{\text{消去旧曲率信息}}
\]
DFP 校的是 $H$（Hessian \textbf{逆}的近似）。方向取 $d^{(k)} = -H_k \nabla f^{(k)}$。
\end{defbox}

\medskip
\refslide{10使用导数的优化算法}{91}{10.52}
\medskip
\refslide{10使用导数的优化算法}{95}{10.56}
\medskip

\begin{notebox}
\textbf{公式记忆}：两个分数项。第一项分子 $s s^T$（秩1矩阵）用位移差修正；第二项用 $Hy$ 消去旧方向上的曲率信息。两项合起来是\emph{秩2校正}。DFP 校 $H$（逆），BFGS 校 $B$（直接近似 Hessian）。
\end{notebox}

\subsubsection{BFGS 简介}

BFGS 是 DFP 的"互补"方法。BFGS 公式：
\[
B_k = B_{k-1} + \frac{y_{k-1} y_{k-1}^T}{y_{k-1}^T s_{k-1}} - \frac{B_{k-1} s_{k-1} s_{k-1}^T B_{k-1}}{s_{k-1}^T B_{k-1} s_{k-1}}
\]
其中 $B_k \approx \nabla^2 f(x^{(k)})$ 是 Hessian 本身的近似。实际使用中，BFGS 往往比 DFP 更稳定（对一维搜索误差不敏感）。考试中以 DFP 为主。

\subsection{8.6 真题示例：第四题——DFP 拟牛顿法完整求解}

\begin{notebox}
\textbf{【2025 真题·第四题】}
\[
\min\ f(X) = x_1^2 + 2x_2^2 + x_1,\quad X^{(0)} = (2,1)^T,\ H_0 = I
\]
用 DFP 方法求解。
\end{notebox}

\medskip

\subsubsection{第 1 轮迭代}

\textbf{梯度}：$\nabla f = (2x_1+1,\ 4x_2)^T$。在 $X^{(0)}$：$\nabla f^{(0)} = (5,4)^T$。

\textbf{方向}：$d^{(0)} = -H_0 \nabla f^{(0)} = -(5,4)^T$（第一步同最速下降）。

\textbf{一维搜索}：$X^{(0)} + \lambda d^{(0)} = (2-5\lambda,\ 1-4\lambda)^T$。
\begin{align*}
\phi(\lambda) &= (2-5\lambda)^2 + 2(1-4\lambda)^2 + (2-5\lambda) \\
&= 25\lambda^2 - 20\lambda + 4 + 2(16\lambda^2 - 8\lambda + 1) + 2 - 5\lambda \\
&= (25+32)\lambda^2 + (-20-16-5)\lambda + (4+2+2) \\
&= 57\lambda^2 - 41\lambda + 8
\end{align*}
令 $\phi'(\lambda) = 114\lambda - 41 = 0$，得 $\lambda_0 = 41/114$。

\[
X^{(1)} = X^{(0)} + \lambda_0 d^{(0)} = \frac{1}{114}(23,-50)^T
\]
\[
\nabla f^{(1)} = \left(\frac{160}{114},\ \frac{-200}{114}\right)^T = \frac{1}{114}(160,-200)^T
\]

\medskip
\refslide{10使用导数的优化算法}{92}{10.53}
\medskip

\subsubsection{DFP 校正}

$s_0 = X^{(1)} - X^{(0)} = -\frac{41}{114}(5,4)^T$。\\
$y_0 = \nabla f^{(1)} - \nabla f^{(0)} = \frac{1}{114}(160-570,\ -200-456)^T = \frac{1}{114}(-410,\ -656)^T = -\frac{41}{114}(10,16)^T$。

计算 DFP 各项：
\begin{align*}
s_0^T y_0 &= \left(-\frac{41}{114}\right)^2 (5,4)\cdot(10,16)
           = \frac{41^2}{114} \\[4pt]
s_0 s_0^T &= \frac{41^2}{114^2}\begin{pmatrix}5\\4\end{pmatrix}
            \begin{pmatrix}5&4\end{pmatrix}
           = \frac{41^2}{114^2}\begin{pmatrix}25&20\\20&16\end{pmatrix} \\[4pt]
y_0^T H_0 y_0 &= \left(\frac{41}{114}\right)^2 (10^2+16^2)
               = \frac{41^2}{114^2}\cdot 356
\end{align*}

代入 DFP 公式得 $H_1$。此处公式展开冗长，核心是：$H_1$ 不再是对角矩阵，它包含了曲率信息。完整结果见课件。

\medskip
\refslide{10使用导数的优化算法}{96}{10.57}
\medskip

\subsubsection{第 2 轮迭代}

$d^{(1)} = -H_1 \nabla f^{(1)}$。一维搜索得 $\lambda_1 = 5/57$。$X^{(2)} = (-1/2,\ 0)^T$。

$\nabla f(X^{(2)}) = (0,0)^T$ → 驻点。$\nabla^2 f = \begin{pmatrix}2&0\\0&4\end{pmatrix}$ 正定 → $X^{(2)}$ 是极小点。

\begin{notebox}
\textbf{DFP 的二次终止性}：对正定二次函数，DFP 法至多 $n$ 轮收敛。此题 $n=2$，恰好两轮到达精确解。这是拟牛顿法的关键理论性质。
\end{notebox}

\medskip
\refslide{10使用导数的优化算法}{102}{10.63}
\medskip

\subsection{本章速查}

\begin{itemize}
  \item 下降算法格式：$x^{(k+1)} = x^{(k)} + \lambda_k d^{(k)}$，$\nabla f^T d < 0$
  \item 共轭梯度法（FR）：$\beta_k = \|\nabla f^{(k)}\|^2 / \|\nabla f^{(k-1)}\|^2$，$d^{(k)} = -\nabla f^{(k)} + \beta_k d^{(k-1)}$
  \item 拟牛顿条件：$H_k y_{k-1} = s_{k-1}$
  \item DFP 公式（秩2校正 $H$）：两项分数修正
  \item BFGS 公式（秩2校正 $B$）：DFP 的互补公式
  \item 二次终止性：对正定二次函数至多 $n$ 步收敛
\end{itemize}

% ============================================================
%  第9章：约束优化方法
% ============================================================
\section{第9章 \quad 约束优化方法}

\subsection{9.1 约束优化方法的两大流派}

第8章的方法适用于\emph{无约束}问题。加上约束后，有两类思路：

\begin{itemize}
  \item \textbf{可行方向法}：始终保持在可行域内，从可行点出发沿下降可行方向移动。每次迭代先找方向（不能跑出可行域），再定步长（确保不越界）。代表：Zoutendijk 法、Rosen 梯度投影法。
  \item \textbf{罚函数法}：将约束问题\emph{转化为一系列无约束问题}。对违反约束的行为加以惩罚，惩罚力度越来越大，迫使无约束问题的解趋向原约束问题的解。代表：外点法（SUMT）、内点法（障碍法）、乘子法。
\end{itemize}

\subsection{9.2 可行方向法通用框架}

假设当前在可行点 $x^{(k)}$。我们需要：
\begin{enumerate}
  \item 找\emph{下降方向} $d$：$\nabla f(x^{(k)})^T d < 0$
  \item 同时 $d$ 是\emph{可行方向}：对充分小的 $\lambda > 0$，$x^{(k)} + \lambda d$ 仍在可行域内
\end{enumerate}

\emph{积极约束}（active constraints）——当前取等号的不等式约束——是限制可行方向的关键。非积极约束在局部可以忽略。

\medskip
\refslide{12可行方向法}{2}{12.1}
\medskip
\refslide{12可行方向法}{3}{12.2}
\medskip

\subsection{9.3 Rosen 梯度投影法}

Rosen 法的核心思想：将\emph{负梯度方向投影到约束的切空间上}，得到既下降又（局部）可行的方向。

\begin{defbox}
\textbf{投影矩阵}

设积极约束的系数矩阵为 $M$（每行是一个积极约束的梯度 $\nabla g_i(x)^T$，仅对线性约束）。则投影矩阵为：
\[
P = I - M^T (M M^T)^{-1} M
\]
$P$ 把任意向量投影到约束的\emph{切空间}（null space of $M$）上。

搜索方向：$d = -P \nabla f$（负梯度在切线空间上的投影）。
\end{defbox}

\medskip
\refslide{12可行方向法}{42}{12.36}
\medskip

\begin{notebox}
\textbf{投影矩阵的几何直觉}：$P \nabla f = 0$ 意味着梯度完全落在约束法向量的张成空间中——这是 KKT 条件的等价表述。$P \nabla f \neq 0$ 意味着梯度在切线方向上有分量，沿此分量移动可以下降而不立即违反约束。
\end{notebox}

\subsubsection{积极约束集的维护}

Rosen 法的难点在于\emph{何时加入约束、何时移除约束}：
\begin{itemize}
  \item 若步长被某个新的不等式约束限制，将其加入积极集
  \item 若 $d = 0$（无法在当前积极集内下降），检查 KKT 乘子符号。若有 $\lambda_i < 0$，表示对应的约束不必是积极的——将其从积极集中移除，重新计算 $P$
\end{itemize}

\subsection{9.4 Zoutendijk 可行方向法}

Zoutendijk 法用线性规划子问题来找方向：在满足所有下降和可行方向约束的条件下，选择使目标函数下降最快的方向。这是求解 LP 子问题来确定搜索方向。

\subsection{9.5 真题示例：第五题——Rosen 梯度投影法}

\begin{notebox}
\textbf{【2025 真题·第五题】}
\[
\begin{aligned}
\min\ & x_1^2 + 2x_2^2 - x_1 x_2 + x_3^2 + x_3 \\
\text{s.t.}\ & x_1 + x_2 = 6 \\
& -x_1 + x_2 + x_3 \le 2 \\
& x_1, x_2, x_3 \ge 0
\end{aligned}
\]
\end{notebox}

\medskip

\subsubsection{化简与初始点}

观察目标函数：$\frac{\partial f}{\partial x_3} = 2x_3 + 1$。在可行域内 $x_3 \ge 0$，故 $x_3 = 0$ 处 $f$ 随 $x_3$ 增大而增大，最优解处 $x_3^* = 0$。问题化为二维。

取初始可行点 $X^{(0)} = (2,2,2)^T$。此时前两个约束取等号（积极）。

\[
\nabla f = (2x_1-x_2,\ 4x_2-x_1,\ 2x_3+1)^T,\quad
\nabla f^{(0)} = (2,\ 6,\ 5)^T
\]

积极约束系数矩阵（取等号的行）：
\[
M = \begin{pmatrix} 1 & 1 & 0 \\ -1 & 1 & 0 \end{pmatrix}
\]
（注意：$x_3 \ge 0$ 即 $-x_3 \le 0$，但 $x_3=2>0$，非积极，暂不列入 $M$。）

\subsubsection{计算投影矩阵}

$MM^T = \begin{pmatrix}2&0\\0&2\end{pmatrix} = 2I$，$(MM^T)^{-1} = \frac{1}{2}I$。
\[
M^T(MM^T)^{-1}M = \frac{1}{2}M^T M = \frac{1}{2}
\begin{pmatrix} 1&-1\\1&1\\0&0 \end{pmatrix}
\begin{pmatrix} 1&1&0\\-1&1&0 \end{pmatrix}
= \frac{1}{2}\begin{pmatrix}2&0&0\\0&2&0\\0&0&0\end{pmatrix}
= \begin{pmatrix}1&0&0\\0&1&0\\0&0&0\end{pmatrix}
\]
\[
P = I - M^T(MM^T)^{-1}M = \begin{pmatrix}0.5&0&-0.5\\0&0&0\\-0.5&0&0.5\end{pmatrix}
\]

搜索方向：$d^{(0)} = -P\nabla f^{(0)} = (1.5,\ 0,\ -1.5)^T$。

\medskip
\refslide{12可行方向法}{44}{12.38}
\medskip

\subsubsection{一维搜索}

$X^{(0)} + \lambda d^{(0)} = (2+1.5\lambda,\ 2,\ 2-1.5\lambda)^T$。代入 $f$ 得单变量函数：
\[
\phi(\lambda) = 4.5\lambda^2 - 4.5\lambda + 14
\]
$\phi'(\lambda) = 9\lambda - 4.5 = 0 \Rightarrow \lambda_0 = 0.5$。

$X^{(1)} = (2.75,\ 2,\ 1.25)^T$。

\subsubsection{第二次迭代与 KKT 判定}

在 $X^{(1)}$ 继续 Rosen 法。新的 $d^{(1)} = 0$（因为 $\nabla f^{(1)}$ 在各积极约束的法向量张成空间中），检查 KKT 乘子。若乘子均非负，则 $X^{(1)}$ 是 KKT 点。

检查 $f$ 的凸性：$\nabla^2 f = \begin{pmatrix}2&-1&0\\-1&4&0\\0&0&2\end{pmatrix}$，顺序主子式 $2>0, 8-1=7>0, 2\times 7 = 14>0$，正定。$f$ 是凸函数，约束为线性 → 凸规划 → KKT 点即全局最优。

\medskip
\refslide{12可行方向法}{46}{12.40}
\medskip

\subsection{9.6 罚函数法}

罚函数法的基本思想：\emph{把约束优化转化为无约束优化}。

\subsubsection{外点法（SUMT——序列无约束极小化技术）}

\begin{defbox}
\textbf{外点罚函数}

对约束问题 $\min\{f(x) \mid g_i(x) \le 0,\ h_j(x)=0\}$，构造罚函数：
\[
P(x, \sigma) = f(x) + \sigma \left(\sum_i [\max\{0, g_i(x)\}]^2 + \sum_j h_j(x)^2 \right)
\]
其中 $\sigma > 0$ 是\emph{罚因子}。当 $x$ 违反约束时，惩罚项为正。$\sigma$ 逐步增大，迫使惩罚项趋近于零。

序列：选递增序列 $\sigma_1 < \sigma_2 < \cdots \to \infty$，对每个 $\sigma_k$ 求解 $\min_x P(x, \sigma_k)$，得 $x^{(k)}$。$x^{(k)} \to x^*$（原问题最优解）。
\end{defbox}

\medskip
\refslide{13罚函数法}{2}{13.1}
\medskip
\refslide{13罚函数法}{3}{13.2}
\medskip

\begin{notebox}
\textbf{为什么叫"外点法"？} 中间迭代点可以\emph{不满足约束}（在可行域外部），随着 $\sigma \to \infty$ 才逐渐趋向可行域。这是"从可行域外部逼近"的意思。
\end{notebox}

\subsubsection{内点法（障碍法）}

与外点法相反，内点法\emph{始终保持在可行域内部}，通过加一个"障碍项"来防止迭代点靠近边界。

\begin{defbox}
\textbf{内点障碍函数}（对仅含不等式约束的问题）

\[
B(x, r) = f(x) + r \sum_i \left(-\frac{1}{g_i(x)}\right) \quad\text{或}\quad f(x) - r \sum_i \ln(-g_i(x))
\]
其中 $r > 0$ 是障碍因子。$r \to 0^+$ 时障碍项的影响逐渐消失。

要求：$x$ 必须严格可行（$g_i(x) < 0$），否则障碍项无定义。
\end{defbox}

\medskip
\refslide{13罚函数法}{15}{13.14}
\medskip

\subsubsection{乘子法（增广 Lagrange 法）}

外点法的缺点是 $\sigma \to \infty$ 时 Hessian 条件数爆炸，数值困难。乘子法通过引入 Lagrange 乘子和一个二次惩罚项，避免了 $\sigma \to \infty$ 的需求。

\begin{defbox}
\textbf{增广 Lagrange 函数}（对等式约束）

\[
L_A(x, \lambda, \sigma) = f(x) + \sum_j \lambda_j h_j(x) + \frac{\sigma}{2}\sum_j h_j(x)^2
\]
乘子更新公式：$\lambda_j^{(k+1)} = \lambda_j^{(k)} + \sigma h_j(x^{(k)})$。

乘子法的关键优势：$\sigma$ 不需要趋于无穷。通过不断更新乘子估计，$\lambda^{(k)}$ 收敛到真实的 Lagrange 乘子，$x^{(k)}$ 收敛到 $x^*$。
\end{defbox}

\medskip
\refslide{13罚函数法}{25}{13.24}
\medskip
\refslide{13罚函数法}{26}{13.25}
\medskip

\subsection{9.7 三种方法的对比}

\begin{defbox}
\textbf{表 9.1：约束优化方法对比}
\[
\begin{array}{l|c|c|c}
& \text{外点法} & \text{内点法} & \text{乘子法} \\ \hline
\text{迭代点位置} & \text{可行域外} & \text{严格可行域内} & \text{均可} \\
\text{参数趋势} & \sigma \to \infty & r \to 0^+ & \sigma\text{ 固定} \\
\text{数值稳定性} & \text{差（条件数大）} & \text{中等} & \text{好} \\
\text{适用性} & \text{等式+不等式} & \text{仅不等式} & \text{等式+不等式}
\end{array}
\]
\end{defbox}

\medskip
\refslide{13罚函数法}{31}{13.30}
\medskip

\subsection{本章速查}

\begin{itemize}
  \item Rosen 投影法：$P = I - M^T(MM^T)^{-1}M$，$d = -P\nabla f$
  \item 积极约束集维护：步长被新约束限 $\to$ 加入；$d=0$ 且有 $\lambda_i<0$ $\to$ 移除
  \item 外点法：$P(x,\sigma) = f(x) + \sigma\sum[\max(0,g_i)]^2$，$\sigma \to \infty$
  \item 内点法：$B(x,r) = f(x) + r\sum(-1/g_i)$，$r \to 0^+$
  \item 乘子法：增广 Lagrange 函数 + 乘子迭代，$\sigma$ 不需趋于无穷
\end{itemize}


\end{document}
